\frfilename{mt286.tex} \versiondate{30.3.16} \copyrightdate{2000}
\newdimen\sparewidth

\def\chaptername{Analisis Fourier} \def\sectionname{Teorema Carleson}

\def\energy{\mathop{\text{energi}}\nolimits}
\def\Innerprod#1#2{\bigl(#1\bigr|#2\bigr)}
\def\mass{\mathop{\text{massa}}\nolimits}
\def\recheck{\discrversionA{\immediate\write0{query}
\global\advance\footnotenumber by 1
\oldfootnote{$^{\the\footnotenumber}$}{recheck}}{}}
\def\vartildef{\tilde{\hbox{$f$}}}

%try:  f\in\eusm L^2
%      h rapidly decr test fn
%      g other

\newsection{286}

Teorema Carleson ({\smc Carleson 66}) merupakan solusi %Dorothy Maharam once said
yang tidak terduga bagi sebuah masalah lama. Hebatnya, teorema ini dapat
dibuktikan dengan metode-metode `elementer'. Bagian tersulit dari uraian di
bawah ini, pada 286J-286L, hanya menuntut pemeriksaan sejumlah
ketaksamaan yang melelahkan. Bagaimana ketaksamaan-ketaksamaan itu dipilih
merupakan persoalan lain; untuk sekali ini, beberapa gagasan dalam pembuktian
terwujud dalam pernyataan lema-lemanya. Argumen di sini merupakan versi yang
sangat diperluas dari {\smc Lacey \& Thiele 00}.

Teorema Maksimal Hardy-Littlewood (286A) penting dan layak dipelajari,
sekalipun Anda membiarkan bagian selebihnya menjadi monumen yang belum
ditelaah. Saya menempatkan 286B-286D %286B 286C 286D
di awal bagian, meskipun ketiganya tidak jauh lebih dari latihan yang
dikerjakan, sebab hasil-hasil itu juga berpotensi berguna dalam konteks lain.

%mention subsupersubscripts?

\cmmnt{Kompleksitas argumen ini membuat kita perlu memperkenalkan cukup
banyak notasi khusus. Alih-alih memasukkan semuanya ke dalam indeks umum,
saya mencantumkannya dalam daftar pada 286W. Di antaranya terdapat sepuluh
konstanta $C_1,\ldots,C_{10}$. Nilai bilangan-bilangan ini tidak penting.
Metode pembuktian di sini sama sekali tidak sesuai jika kita hendak menaksir
laju konvergensi. Saya memberikan cara menghitung $C_n$ semata-mata demi
alur logika linear yang digunakan dalam risalah ini, dan karena cara tersebut
sesekali memberi petunjuk tentang taktik yang sedang digunakan. }

Dalam bagian ini semua integral diambil terhadap ukuran Lebesgue $\mu$ pada
$\Bbb R$ kecuali dinyatakan sebaliknya.

\leader{286A}{Teorema Maksimal} % Hardy \& Littlewood 30, but not best constant
 Misalkan $1<p<\infty$ dan $f\in\eusm L_{\Bbb C}^p(\mu)$\cmmnt{
(definisi: 244P)}. Tetapkan

\Centerline{$f^*(x)=\sup\{\Bover1{b-a}\int_a^b|f|:a\le x\le b$, $a<b\}$}

\noindent untuk $x\in\Bbb R$. Maka
$\|f^*\|_p\le\Bover{2^{1/p}p}{p-1}\|f\|_p$.

\proof{{\bf (a)} Cukup kita tinjau kasus $f=|f|$. Perhatikan bahwa jika
$E\subseteq\Bbb R$ mempunyai ukuran berhingga, maka

\Centerline{$\int_Ef =\int(f\times\chi E)\times\chi E \le\|f\times\chi
E\|_p(\mu E)^{1/q}\le\|f\|_p(\mu E)^{1/q}$}

\noindent berhingga, dengan $q=\Bover{p}{p-1}$, menurut ketaksamaan
H\"older (244Eb). Akibatnya, jika $t>0$ dan $\int_Ef\ge t\mu E$, haruslah
$t\mu E\le\|f\times\chi E\|_p(\mu E)^{1/q}$, $t(\mu
E)^{1/p}\le\|f\times\chi E\|_p$ dan

\Centerline{$\mu E\le\Bover1{t^p}\|f\times\chi E\|^p_p
=\Bover1{t^p}\int_Ef^p$.}

\medskip

{\bf (b)} Untuk $t>0$, tetapkan

\Centerline{$G_t=\{x:t(y-x)<\int_x^yf$ untuk beberapa $y>x\}$.}

\medskip

\quad{\bf (i)} $G_t$ merupakan himpunan terbuka. \Prf\ Untuk setiap $y\in\Bbb R$,

\Centerline{$G_{ty}=\{x:x<y$, $t(y-x)<\int_x^yf\}$}

\noindent terbuka, karena $x\mapsto t(y-x)$ dan $x\mapsto\int_x^yf$
kontinu (225A); jadi $G_t=\bigcup_{y\in\Bbb R}G_{ty}$ terbuka.\ \Qed

\medskip

\quad{\bf (ii) } Menurut 2A2I, terdapat partisi $\Cal C$ atas $G_t$ menjadi
interval-interval terbuka. Sekarang $C$ terbatas dan
$t\mu C\le\int_Cf$ untuk setiap
$C\in\Cal C$.

\medskip

\quad\Prf\grheada\ Untuk $x\in C$, tinjau $F_x=\{y:y\ge x$,
$t(y-x)\le\int_x^yf\}$. $x\in F_x$ dan
$y-x\le\Bover1{t^p}\int_{-\infty}^{\infty}f^p$ untuk setiap $y\in F_x$, dengan
(a), sehingga $F_x$ terbatas di atas. Tetapkan $z_x=\sup F_x$. Karena
$y\mapsto t(y-x)-\int_x^yf$ adalah kontinu, $z_x\in F_x$. \Quer\ Jika $z_x\in
G_t$, terdapat $y>z_x$ sedemikian sehingga
$t(y-z_x)<\int_{z_x}^yf$; tetapi sekarang

\Centerline{$t(y-x) \le\int_x^{z_x}f+\int_{z_x}^yf=\int_x^yf$}

\noindent dan $y\in F_x$, suatu kemustahilan.\ \BanG\ Jadi
$z_x\notin G_t$ dan $z_x\notin C$, sehingga $z_x$ merupakan batas atas $C$.

\medskip

\qquad\grheadb\ Ini menunjukkan bahwa

\Centerline{$\sup C\le z_x\le x+\Bover1{t^p}\int_{-\infty}^{\infty}f^p$}

\noindent untuk setiap $x\in C$. Jadi sebenarnya $C$ terbatas dan berbentuk
$\ooint{a,b}$ dengan $a<b$ dalam $\Bbb R$. \Quer\ Jika
$t(b-a)>\int_a^bf$, ada sebuah $x\in\ooint{a,b}$ sedemikian sehingga
$t(b-x)>\int_x^bf$. Kita mengetahui bahwa $b\le z_x$ dan $b\notin G_t$,
sehingga $t(z_x-b)\ge\int_b^{z_x}f$. Dengan menjumlahkan kedua ketaksamaan,
$t(z_x-x)>\int_x^{z_x}f$ dan
$z_x\notin F_x$.\ \Bang

\medskip

\qquad\grheadc\ Jadi $t\mu C\le\int_Cf$, sebagaimana dinyatakan.\ \Qed

\medskip

\quad{\bf (iii)} Oleh karena itu, karena $\Cal C$ terhitung dan $f$
nonnegatif, kita dapat menerapkan (a) dalam bentuk penuhnya untuk
melihat bahwa

\Centerline{$\mu G_t=\sum_{C\in\Cal C}\mu C \le\sum_{C\in\Cal
C}\Bover1{t^p}\int_Cf^p \le\Bover1{t^p}\int_{-\infty}^{\infty}f^p$}

\noindent adalah hingga, dan

\Centerline{$\int_{G_t}f=\sum_{C\in\Cal C}\int_Cf \ge\sum_{C\in\Cal C}t\mu
C=t\mu G_t$.}

\medskip

{\bf (c)} Semua ini berlaku untuk setiap $t>0$. Sekarang tetapkan

\Centerline{$f_1^*(x)=\sup_{b>x}\Bover1{b-x}\int_x^bf$}

\noindent untuk $x\in\Bbb R$, kita memiliki $\{x:f_1^*(x)>t\}=G_t$ untuk
setiap $t>0$.

Untuk setiap $t>0$,

\Centerline{$\Bover1pt\mu G_t =(1-\Bover1q)t\mu G_t \le\int_{G_t}f-\Bover1q
t\chi\Bbb R \le\int_{-\infty}^{\infty}(f-\Bover1q t\chi\Bbb R)^+$.}

\noindent Jadi

$$\eqalignno{\int_{-\infty}^{\infty}(f_1^*)^p
&=\int_0^{\infty}\mu\{x:f_1^*(x)^p>t\}dt\cr
\displaycause{lihat 252O}
&=p\int_0^{\infty}u^{p-1}\mu\{x:f_1^*(x)>u\}du\cr
\displaycause{dengan substitusi $t=u^p$}
&=p\int_0^{\infty}u^{p-1}\mu G_udu
\le p^2\int_0^{\infty}u^{p-2}
  \bigl(\int_{-\infty}^{\infty}(f-\Bover1q u\chi\Bbb R)^+\bigr)du\cr
&=p^2\int_{-\infty}^{\infty}\int_0^{\infty}
  \max(0,f(x)-\Bover1q u)u^{p-2}dudx\cr
\displaycause{menurut teorema Fubini, 252B, karena
$(x,u)\mapsto u^{p-2}\max(0,f(x)-\bover1qu)$ terukur dan
nonnegatif}
&=p^2\int_{-\infty}^{\infty}\int_0^{qf(x)}
  u^{p-2}(f(x)-\Bover1q u)dudx\cr
&=\Bover{p^2q^{p-1}}{p(p-1)}\int_{-\infty}^{\infty}f^p
=(\Bover{p}{p-1})^p\|f\|_p^p.\cr}$$

\medskip

{\bf (d)} Demikian pula, dengan menetapkan
$f_2^*(x)=\sup_{a<x}\Bover1{x-a}\int_a^xf$
untuk $x\in\Bbb R$,
$\int_{-\infty}^{\infty}(f_2^*)^p\le(\Bover{p}{p-1})^p\|f\|_p^p$. Tetapi
$f^*=\max(f_1^*,f_2^*)$. \Prf\ Tentu saja $f_1^*\le f^*$ dan $f_2^*\le f^*$.
Selain itu, jika $f^*(x)>t$, harus ada interval tak trivial $I$ yang memuat
$x$ sedemikian sehingga $\int_If>t\mu I$; jika $a=\inf I$ dan
$b=\sup I$, maka
$\int_a^xf>(x-a)t$ dan $f_2^*(x)>t$, atau $\int_x^bf>(b-x)t$ dan $f_1^*(x)>t$.
Karena $x$ dan $t$ sebarang, $f^*=\max(f_1^*,f_2^*)$.\ \Qed

Oleh karena itu

$$\eqalign{\|f^*\|_p^p
&=\int_{-\infty}^{\infty}(f^*)^p
=\int_{-\infty}^{\infty}\max((f_1^*)^p,(f_2^*)^p)\cr
&\le\int_{-\infty}^{\infty}(f_1^*)^p+(f_2^*)^p
\le 2(\Bover{p}{p-1})^p\|f\|^p_p.\cr}$$

\noindent Dengan mengambil akar pangkat $p$, kita memperoleh ketaksamaan
yang dicari. }%end of proof of 286A

\leader{286B}{Lema} Misalkan $g:\Bbb R\to\coint{0,\infty}$ suatu fungsi yang
tidak menurun pada $\ocint{-\infty,\alpha}$, tidak meningkat pada
$\coint{\beta,\infty}$ dan konstan pada $[\alpha,\beta]$, dengan
$\alpha\le\beta$. Maka untuk setiap fungsi terukur $f:\Bbb
R\to[0,\infty]$, $\int_{-\infty}^{\infty}f\times g\le\int_{-\infty}^{\infty}g
\cdot\sup_{a\le\alpha,b\ge\beta,a<b}\Bover1{b-a}\int_a^bf$.

\proof{ Tetapkan
$\gamma=\sup_{a\le\alpha,b\ge\beta,a<b}\Bover1{b-a}\int_a^bf$.
Untuk $n$, $k\in\Bbb N$, tetapkan
$E_{nk}=\{x:\alpha-2^n\le x\le\beta+2^n$, $g(x)\ge
2^{-n}(k+1)\}$. Maka $E_{nk}$ kosong atau merupakan interval terbatas yang
memuat $[\alpha,\beta]$, dan $\int_{E_{nk}}f\le\gamma\mu E_{nk}$. Untuk
$n\in\Bbb N$, tetapkan $g_n=2^{-n}\sum_{k=0}^{4^n-1}\chi E_{nk}$; maka
$\sequencen{g_n}$ merupakan barisan fungsi tak menurun dengan supremum $g$,
dan

$$\eqalign{\int_{-\infty}^{\infty} f\times g
&=\sup_{n\in\Bbb N}\int_{-\infty}^{\infty} f\times g_n
=\sup_{n\in\Bbb N}2^{-n}\sum_{k=0}^{4^n-1}\int_{E_{nk}}f\cr
&\le\sup_{n\in\Bbb N}2^{-n}\sum_{k=0}^{4^n-1}\gamma\mu E_{nk}
=\sup_{n\in\Bbb N}\gamma\int_{-\infty}^{\infty} g_n
=\gamma\int_{-\infty}^{\infty} g,\cr}$$

\noindent sebagaimana dinyatakan. }%end of proof of 286B

\cmmnt{\medskip

\noindent{\bf Catatan} Bandingkan dengan 224J.}

\vleader{48pt}{286C}{Pergeseran, modulasi, dan dilasi}\cmmnt{ Beberapa
perhitungan di bawah akan lebih mudah jika kita menggunakan formalisme
berikut.} Untuk setiap fungsi $f$ yang domainnya termuat dalam $\Bbb R$,
dan $\alpha\in\Bbb R$, kita dapat mendefinisikan

\Centerline{$(S_{\alpha}f)(x)=f(x+\alpha)$, \quad$(M_{\alpha}f)(x)=e^{i\alpha
x}f(x)$, \quad$(D_{\alpha}f)(x)=f(\alpha x)$}

\noindent kapan pun ruas-ruas kanannya terdefinisi. \cmmnt{ Dalam hal
$S_{\alpha}f$ dan $D_{\alpha}f$ kadang-kadang nyaman untuk membiarkan
$\pm\infty$ sebagai nilai fungsi tersebut. Kita mempunyai fakta-fakta
elementer berikut.}

\spheader 286Ca $S_{-\alpha}S_{\alpha}f=f$, $D_{1/\alpha}D_{\alpha}f=f$ jika
$\alpha\ne 0$.

\spheader 286Cb $S_{\alpha}(f\times g)=S_{\alpha}f\times S_{\alpha}g$,
$D_{\alpha}(f\times g)=D_{\alpha}f\times D_{\alpha}g$.

\spheader 286Cc $D_{\alpha}|f|=|D_{\alpha}f|$.

\spheader 286Cd Jika $f$ terintegralkan, maka

\Centerline{$(M_{\alpha}f)\varsphat=S_{-\alpha}\varhatf$,
\quad$(S_{\alpha}f)\varsphat=M_{\alpha}\varhatf$,
\quad$(S_{\alpha}f)\varspcheck=M_{-\alpha}\varcheckf$;}

\noindent jika selain itu $\alpha>0$, maka

\Centerline{$\alpha(D_{\alpha}f)\varsphat =D_{1/\alpha}\varhatf$,
\quad$\alpha(D_{\alpha}f)\varspcheck =D_{1/\alpha}\varcheckf$\dvro{.}{}}

\cmmnt{\noindent (283Cc-283Ce).}

\spheader 286Ce Jika $f$ termasuk dalam
$\eusm L^1_{\Bbb C}=\eusm L^1_{\Bbb C}(\mu)$, demikian pula
$S_{\alpha}f$, $M_{\alpha}f$, dan (jika $\alpha\ne 0$) $D_{\alpha}f$;
dalam hal ini

\Centerline{$\|S_{\alpha}f\|_1=\|M_{\alpha}f\|_1=\|f\|_1$,
\quad$\|D_{\alpha}f\|_1=\Bover1{|\alpha|}\|f\|_1$.}

\spheader 286Cf Jika $f$ termasuk dalam $\eusm L^2_{\Bbb C}$, demikian pula
$S_{\alpha}f$, $M_{\alpha}f$, dan (jika $\alpha\ne 0$) $D_{\alpha}f$;
dalam hal ini

\Centerline{$\|S_{\alpha}f\|_2=\|M_{\alpha}f\|_2=\|f\|_2$,
\quad$\|D_{\alpha}f\|_2=\Bover1{\sqrt{|\alpha|}}\|f\|_2$.}

\spheader 286Cg Jika $h$ merupakan fungsi uji yang menurun cepat\cmmnt{
(284A)}, demikian pula $M_{\alpha}h$, $S_{\alpha}h$, dan
(jika $\alpha\ne 0$) $D_{\alpha}h$.

\leader{286D}{Lema} Misalkan $g:\Bbb R\to[0,\infty]$ suatu fungsi terukur
sedemikian sehingga, untuk suatu konstanta $C\ge 0$,
$\int_Eg\le C\sqrt{\mu E}$ kapan pun $\mu E<\infty$. Maka $g$ berhingga
hampir di mana-mana dan
$\biggerint_{-\infty}^{\infty}\Bover1{1+|x|}g(x)dx$ berhingga.

\proof{ Untuk setiap $n\ge 1$, tetapkan
$E_n=\{x:|x|\le n,\,g(x)\ge n\}$; maka

\Centerline{$n\mu E_n\le\int_{E_n}g\le C\sqrt{\mu E_n}$,}

\noindent jadi $\mu E_n\le\Bover{C^2}{n^2}$ dan

\Centerline{$\{x:g(x)=\infty\}=\bigcap_{n\ge 1}\bigcup_{m\ge n}E_m$}

\noindent mempunyai ukuran paling besar
$\inf_{n\ge 1}\sum_{m=n}^{\infty}\mu E_m=0$.

Untuk integralnya, tetapkan $G(x)=\int_0^xg$ untuk $x\ge 0$. Maka, untuk setiap
$a\ge 0$,

$$\eqalignno{\int_0^a\Bover{g(x)}{1+x}dx
&=\Bover{G(a)}{1+a}+\int_0^a\Bover{G(x)}{(1+x)^2}dx\cr
\displaycause{225F}
&\le C\bigl(\Bover{\sqrt{a}}{1+a}
  +\int_0^a\Bover{\sqrt{x}}{(1+x)^2}dx\bigr)
\le C\bigl(1+\int_0^{\infty}\Bover{\sqrt{x}}{(1+x)^2}dx\bigr),\cr}$$

\noindent jadi

\Centerline{$\int_0^{\infty}\Bover{g(x)}{1+x}dx \le
C\bigl(1+\int_0^{\infty}\Bover{\sqrt{x}}{(1+x)^2}dx\bigr)$}


\noindent berhingga. Demikian pula,
$\biggerint_{-\infty}^0\Bover{g(x)}{1-x}dx$ berhingga, sehingga hasilnya
terbukti. }%end of proof of 286D

\allowmorestretch{500}{ \leader{286E}{Konstruksi Lacey-Thiele (a)}
Misalkan $\Cal I$ keluarga semua {\bf interval diadik} berbentuk
\penalty-100$\coint{2^kn,2^k(n+1)}$, dengan $k$, $n\in\Bbb Z$. Sifat
geometris penting $\Cal I$ ialah bahwa, untuk $I$, $J\in\Cal I$, berlaku
$I\subseteq J$, atau $J\subseteq I$, atau $I\cap J=\emptyset$.
Misalkan $Q$ himpunan semua pasangan
$\sigma=(I_{\sigma},J_{\sigma})\in\Cal I^2$ sedemikian sehingga
$\mu I_{\sigma}\cdot\mu J_{\sigma}=1$. Untuk $\sigma\in Q$, misalkan
$k_{\sigma}\in\Bbb Z$ memenuhi $\mu J_{\sigma}=2^{k_{\sigma}}$ dan
$\mu I_{\sigma}=2^{-k_{\sigma}}$; misalkan $x_{\sigma}$ titik tengah
$I_{\sigma}$, $y_{\sigma}$ titik tengah $J_{\sigma}$,
$J^l_{\sigma}\in\Cal I$ paruh kiri $J_{\sigma}$,
$J^r_{\sigma}\in\Cal I$ paruh kanan $J_{\sigma}$, dan $y^l_{\sigma}$
titik seperempat bawah $J_{\sigma}$\cmmnt{, yaitu titik tengah
$J^l_{\sigma}$}. }%end of allowmorestretch

\spheader 286Eb Terdapat fungsi uji yang menurun cepat $\phi$ sedemikian
sehingga $\varhat{\phi}$ bernilai real dan
$\chi[-\bover16,\bover16]\le\varhat{\phi}\le\chi[-\bover15,\bover15]$.
\prooflet{\Prf\ Lihat bagian (b)-(d) pembuktian 284G. Proses di sana dapat
digunakan untuk memperoleh fungsi mulus $\psi_1$ yang nol di luar interval
$[\bover16,\bover15]$ dan positif tegas pada
$\ooint{\bover16,\bover15}$; dengan mengalikannya dengan faktor yang sesuai,
kita dapat mengatur agar $\int_{-\infty}^{\infty}\psi_1=1$. Jadi, jika
ditetapkan $\psi_2(x)=1-\int_{-\infty}^x\psi_1$ untuk $x\in\Bbb R$,
maka $\psi_2$ mulus dan
$\chi\ocint{-\infty,\bover16}\le\psi_2\le\chi\ocint{-\infty,\bover15}$.
Sekarang tetapkan $\psi_0(x)=\psi_2(x)\psi_2(-x)$ untuk $x\in\Bbb R$,
dan $\phi=\varcheck{\psi}_0$; maka $\varhat{\phi}=\psi_0$ (284C)
mempunyai sifat yang diperlukan.\ \Qed} %wouldn't it be nicer if $\|\phi\|_2=1$?  but it is helpful at one
%point to have $\varhat\phi(0)=1$

Untuk $\sigma\in Q$, tetapkan\cmmnt{ $\phi_{\sigma}
=2^{k_{\sigma}/2}M_{y^l_{\sigma}}S_{-x_{\sigma}}D_{2^{k_{\sigma}}}\phi$,
sehingga}

\Centerline{$\phi_{\sigma}(x) =\sqrt{\mu J_{\sigma}}e^{iy^l_{\sigma}x}
\phi((x-x_{\sigma})\mu J_{\sigma})$.}

\noindent\cmmnt{Perhatikan bahwa }$\phi_{\sigma}$ merupakan fungsi uji yang
menurun cepat. \cmmnt{Sekarang
$\varhat{\phi}_{\sigma}=2^{-k_{\sigma}/2}S_{-y^l_{\sigma}}
M_{-x_{\sigma}}D_{2^{-k_{\sigma}}}\varhat{\phi}$, yaitu,}

\Centerline{$\varhat{\phi}_{\sigma}(y) =\sqrt{\mu
I_{\sigma}}e^{-ix_{\sigma}(y-y^l_{\sigma})} \varhat{\phi}((y-y^l_{\sigma})\mu
I_{\sigma})$,}

\noindent yang nol kecuali \cmmnt{ $|y-y^l_{\sigma}|\le\bover15\mu
J_{\sigma}$; karena panjang $J^l_{\sigma}$ adalah $\bover12\mu J_{\sigma}$,
ini hanya mungkin terjadi apabila } $y\in J^l_{\sigma}$. \cmmnt{Kita mempunyai
fakta sederhana berikut.}

\quad(i) $\|\phi_{\sigma}\|_2 =\cmmnt{\sqrt{\mu J_{\sigma}}\cdot\sqrt{\mu
I_{\sigma}}\|\phi\|_2 =}\|\phi\|_2$ untuk setiap $\sigma\in Q$.

\quad(ii) $\|\varhat{\phi}_{\sigma}\|_1 =\cmmnt{\sqrt{\mu I_{\sigma}}\cdot\mu
J_{\sigma}\|\varhat{\phi}\|_1 =}\sqrt{\mu J_{\sigma}}\|\varhat{\phi}\|_1$
untuk setiap $\sigma\in Q$.

\quad(iii) Jika $\sigma$, $\tau\in Q$ dan $J^l_{\sigma}\cap
J^l_{\tau}=\emptyset$ maka

\Centerline{$\innerprod{\phi_{\sigma}}{\phi_{\tau}}
=\innerprod{\varhat{\phi}_{\sigma}}{\varhat{\phi}_{\tau}}=0$\dvro{.}{,}}

\noindent\cmmnt{menurut 284Ob. }(Untuk $f$, $g\in\eusm L_{\Bbb C}^2$, saya menulis
$\innerprod{f}{g}$ untuk $\int_{-\infty}^{\infty}f\times\bar g$.)

\quad(iv) Jika $\sigma$, $\tau\in Q$ dan $J_{\sigma}\ne J_{\tau}$ dan
$J^r_{\sigma}\cap J^r_{\tau}$ tidak kosong, maka \cmmnt{
$J^l_{\sigma}\cap
J^l_{\tau}=\emptyset$ jadi} $\innerprod{\phi_{\sigma}}{\phi_{\tau}}=0$.

\spheader 286Ec Tetapkan $w(x)=\Bover1{(1+|x|)^3}$ untuk $x\in\Bbb R$.
Untuk $\sigma\in Q$, tetapkan\cmmnt{ $w_{\sigma}
=2^{k_{\sigma}}S_{-x_{\sigma}}D_{2^{k_{\sigma}}}w$, sehingga bahwa}

\Centerline{$w_{\sigma}(x)=w((x-x_{\sigma})\mu J_{\sigma})\mu J_{\sigma}
\le\mu J_{\sigma}=2^{k_{\sigma}}$}

\noindent untuk setiap $x$. Perhatikan bahwa $w_{\sigma}=w_{\tau}$ kapan pun
$I_{\sigma}=I_{\tau}$.

\leader{286F}{Urutan parsial (a)} Untuk $\sigma$, $\tau\in Q$, katakan bahwa
$\tau\le\sigma$ jika $J_{\tau}\subseteq J_{\sigma}$ dan
$I_{\sigma}\subseteq I_{\tau}$. Maka $\le$ merupakan urutan parsial pada
$Q$. \cmmnt{Kita mempunyai fakta-fakta elementer berikut. }

\quad(i) Jika $\tau\le\sigma$, maka $k_{\tau}\le k_{\sigma}$.

\quad(ii) Jika $\sigma$ dan $\tau$ tak terbandingkan\cmmnt{ (yaitu,
$\sigma\not\le\tau$ dan $\tau\not\le\sigma$)}, maka $(I_{\sigma}\times
J_{\sigma})\cap(I_{\tau}\times J_{\tau})$ kosong. \prooflet{\Prf\ Kita boleh
menganggap bahwa $k_{\sigma}\le k_{\tau}$. Jika $J_{\sigma}\cap
J_{\tau}\ne\emptyset$, maka $J_{\sigma}\subseteq J_{\tau}$, karena keduanya
merupakan interval diadik dan $J_{\sigma}$ lebih pendek; tetapi karena
$\sigma\not\le\tau$, ini berarti bahwa $I_{\tau}\not\subseteq
I_{\sigma}$ dan $I_{\sigma}\cap I_{\tau}=\emptyset$.\ \Qed}

\quad(iii) Jika $\sigma$, $\sigma'$ tak terbandingkan dan keduanya
lebih besar dari atau sama dengan $\tau$, maka $I_{\sigma}\cap
I_{\sigma'}=\emptyset$\cmmnt{, karena $J_{\tau}\subseteq J_{\sigma}\cap
J_{\sigma'}$}.

\quad(iv) Jika $\tau\le\sigma$ dan $k_{\tau}\le k\le k_{\sigma}$, maka ada
(unik) $\upsilon$ sedemikian sehingga $\tau\le\upsilon\le\sigma$ dan
$k_{\upsilon}=k$. \prooflet{(Intinya ialah bahwa terdapat tepat satu
$I\in\Cal I$ sedemikian sehingga
$I_{\sigma}\subseteq I\subseteq I_{\tau}$ dan $\mu I=2^{-k}$; dan dengan cara
yang sama hanya ada satu kandidat untuk $J_{\upsilon}$.)}

\spheader 286Fb\dvro{ Jika}{ Akan berguna untuk memakai singkatan berikut:
jika} $R\subseteq Q$, tetapkan

\Centerline{$R^+=\bigcup_{\tau\in R}\{\sigma:\tau\le\sigma\in Q\}$.}

 
\spheader 286Fc Untuk $\tau\in Q$, tetapkan

\Centerline{$T_{\tau}=\{\sigma:\sigma\in Q$, $\tau\le\sigma$,
$J^r_{\tau}\subseteq J^r_{\sigma}\}$.}

\noindent Perhatikan bahwa jika $\sigma$, $\sigma'\in T_{\tau}$ dan
$k_{\sigma}\ne k_{\sigma'}$ maka\cmmnt{ $J_{\sigma}\ne J_{\sigma'}$ dan
$J^r_{\sigma}\cap J^r_{\sigma'}\ne\emptyset$, jadi}
$\innerprod{\phi_{\sigma}}{\phi_{\sigma'}}=0$\cmmnt{ (286E(b-iv))}.

\vleader{72pt}{286G}{}\cmmnt{ Kita akan memerlukan hasil beberapa
perhitungan elementer.

\medskip

\noindent}{\bf Lema} (a)
$\int_{-\infty}^{\infty}w_{\sigma}=\int_{-\infty}^{\infty}w=1$ untuk setiap
$\sigma\in Q$.

(b) Untuk setiap $m\in\Bbb N$,
$\sum_{n=m}^{\infty}w(n+\bover12)\le\Bover1{2(1+m)^2}$.

(c) Misalkan $\sigma\in Q$ dan $I$ suatu interval yang tidak memuat
$x_{\sigma}$. Maka $\int_Iw_{\sigma}\ge
w_{\sigma}(x)\mu I$, di mana $x$ adalah titik tengah $I$.

(d) Untuk setiap $x\in\Bbb R$, $\sum_{n=-\infty}^{\infty}w(x-n)\le 2$.

(e) Terdapat konstanta $C_1\ge 0$ sedemikian sehingga $|\phi(x)|\le
C_1\min(w(3),w(x)^2)$ untuk setiap $x\in\Bbb R$ dan

\Centerline{$|\phi_{\sigma}(x)| \le C_1\sqrt{\mu I_{\sigma}}w_{\sigma}(x)
\min(1,w_{\sigma}(x)\mu I_{\sigma})$}

\noindent untuk setiap $x\in\Bbb R$ dan $\sigma\in Q$.

(f) Terdapat konstanta $C_2\ge 0$ sedemikian sehingga
$\int_{-\infty}^{\infty}w(x)w(\alpha x+\beta)dx\le C_2w(\beta)$ kapan pun
$0\le\alpha\le 1$ dan $\beta\in\Bbb R$.

(g) Terdapat konstanta $C_3\ge 0$ sedemikian sehingga
$|\innerprod{\phi_{\sigma}}{\phi_{\tau}}|
\le C_3\sqrt{\mu I_{\sigma}}\sqrt{\mu J_{\tau}}\int_{I_{\tau}}w_{\sigma}$
setiap kali $\sigma$, $\tau\in Q$ dan $k_{\sigma}\le k_{\tau}$.

(h) Terdapat konstanta $C_4\ge 0$ sedemikian sehingga

\Centerline{$\sum_{\sigma\in Q,\sigma\ge\tau,k_{\sigma}=k} \int_{\Bbb
R\setminus I_{\tau}}w_{\sigma}\le C_4$}

\noindent kapan pun $\tau\in Q$ dan $k\in\Bbb Z$.

\proof{{\bf (a)} Langsung dari definisi pada 286Ec, rumus pada 286Ce, dan fakta
bahwa $\biggerint_0^{\infty}\Bover1{(1+x)^3}dx=\Bover12$.

\medskip

{\bf (b)} Intinya ialah bahwa $w$ konveks pada $\ocint{-\infty,0}$
dan $\coint{0,\infty}$. Jadi kita dapat menerapkan 233Ib dengan $f(x)=x$, atau
berdebat langsung dari fakta bahwa
$w(n+\bover12)\le\bover12(w(n+\bover12+x)+w(n+\bover12-x))$ untuk
$|x|\le\bover12$, untuk melihat bahwa $w(n+\bover12)\le\int_n^{n+1}w$ untuk
setiap $n\ge 0$.

\Centerline{$\sum_{n=m}^{\infty}w(n+\bover12)
\le\int_m^{\infty}w=\Bover1{2(1+m)^2}$.}

\medskip

{\bf (c)} Demikian juga, karena $I$ terletak di sisi yang sama dari
$x_{\sigma}$, $w_{\sigma}$ konveks pada $I$, sehingga ketaksamaan
yang sama menghasilkan $w_{\sigma}(x)\mu I\le\int_Iw_{\sigma}$.

\medskip

{\bf (d) } Misalkan $m$ memenuhi $|x-m|\le\bover12$. Kemudian, dengan
menggunakan ketaksamaan yang sama seperti sebelumnya untuk menaksir $w(x-n)$
bagi $n\ne m$, kita memperoleh

$$\eqalign{\sum_{n=-\infty}^{\infty}w(x-n)
&\le w(x-m)+\int_{-\infty}^{x-m-\bover12}w
  +\int_{x-m+\bover12}^{\infty}w\cr
&\le 1+\int_{-\infty}^{\infty}w
=2.\cr}$$

\medskip

{\bf (e)} Karena $\lim_{x\to\infty}x^6\phi(x)=\lim_{x\to-\infty}x^6\phi(x)=0$,
ada $C_1>0$ sedemikian sehingga $|\phi(x)|\le C_1\min(w(3),w(x)^2)$ untuk
setiap $x\in\Bbb R$. Sekarang
$|\phi(x)|\le C_1w(x)^2=C_1w(x)\min(1,w(x))$
untuk setiap $x$, jadi

$$\eqalign{|\phi_{\sigma}(x)|
&=\sqrt{\mu J_{\sigma}}|\phi((x-x_{\sigma})\mu J_{\sigma}|
\le C_1\sqrt{\mu J_{\sigma}}w((x-x_{\sigma})\mu J_{\sigma})
   \min(1,w((x-x_{\sigma})\mu J_{\sigma}))\cr
&=C_1\sqrt{\mu J_{\sigma}}w_{\sigma}(x)\mu I_{\sigma}
   \min(1,w_{\sigma}(x)\mu I_{\sigma})
=C_1\sqrt{\mu I_{\sigma}}w_{\sigma}(x)
   \min(1,w_{\sigma}(x)\mu I_{\sigma})\cr}$$

\noindent setiap kali $\sigma\in Q$ dan $x\in\Bbb R$.

\medskip

{\bf (f)(i)} Pertama, perhatikan bahwa

$$\bover{w(\bover12(1+\beta))}{w(\beta)}
=\bover{8(1+\beta)^3}{(3+\beta)^3}\le 8$$

\noindent untuk setiap $\beta\ge 0$. Sekarang $\alpha w(\alpha+\alpha\beta)\le
4w(\beta)$ kapan pun $\beta\ge 0$ dan $\alpha\ge\bover12$. \Prf\ Untuk
$t\ge\bover12$,

\Centerline{$\Bover{d}{dt}tw(t+t\beta)
=\Bover{1-2t(1+\beta)}{(1+t+t\beta)^4}\le 0$,}

\noindent jadi

\Centerline{$\alpha w(\alpha+\alpha\beta)
\le\bover12w(\bover12+\bover12\beta)\le 4w(\beta)$. \Qed}

Tentu saja ini berarti bahwa

\Centerline{$\Bover1{\alpha}w(\Bover{1+\beta}{2\alpha})\le 8w(\beta)$}

\noindent kapan pun $\beta\ge 0$ dan $0<\alpha\le 1$.

\medskip

\quad{\bf (ii)} Ambil $C_2=16$. Jika $0<\alpha\le 1$ dan $\beta\ge 0$, tetapkan
$\gamma=\Bover{1+\beta}{2\alpha}$. Kemudian, untuk setiap $x\ge-\gamma$,

\Centerline{$1+\alpha x+\beta =(1+\beta)(1+\Bover{\alpha
x}{1+\beta})\ge\Bover12(1+\beta)$,}

\noindent jadi $w(\alpha x+\beta)\le 8w(\beta)$ dan

\Centerline{$\int_{-\gamma}^{\infty}w(x)w(\alpha x+\beta)dx \le
8w(\beta)\int_{-\gamma}^{\infty}w\le 8w(\beta)$.}

\noindent Di sisi lain,

$$\eqalign{\int_{-\infty}^{-\gamma}w(x)w(\alpha x+\beta)dx
&\le w(\gamma)\int_{-\infty}^{\infty}w(\alpha x+\beta)dx\cr
&=\Bover1{\alpha}w(\Bover{1+\beta}{2\alpha})
  \int_{-\infty}^{\infty}w
\le 8w(\beta).\cr}$$

\noindent Menggabungkan ini, $\int_{-\infty}^{\infty}w(x)w(\alpha x+\beta)dx
\le 16w(\beta)$; dan ini benar setiap kali $0<\alpha\le 1$ dan $\beta\ge 0$.

\medskip

\quad{\bf (iii)} Jika $\alpha=0$, maka

\Centerline{$\int_{-\infty}^{\infty}w(x)w(\alpha x+\beta)dx
=w(\beta)\int_{-\infty}^{\infty}w=w(\beta)\le C_2w(\beta)$}

\noindent untuk setiap $\beta$. Jika $0<\alpha\le 1$ dan $\beta<0$, maka

$$\eqalignno{\int_{-\infty}^{\infty}w(x)w(\alpha x+\beta)dx
&=\int_{-\infty}^{\infty}w(-x)w(-\alpha x-\beta)dx\cr
\displaycause{karena $w$ merupakan fungsi genap}
&=\int_{-\infty}^{\infty}w(x)w(\alpha x-\beta)dx
\le C_2w(-\beta)\cr
\displaycause{menurut (ii) di atas}
&=C_2w(\beta).\cr}$$

\noindent Jadi kita memperoleh ketaksamaan yang diperlukan dalam semua
kasus.

\medskip

{\bf (g)} Tetapkan
$C_3=\max(C_1^2C_2,\|\phi\|_2^2/\int_{-1/2}^{1/2}w)$.

\medskip

\quad{\bf (i)} Kasus $\sigma=\tau$ dapat langsung diselesaikan. Dalam hal
ini,

\Centerline{$|\innerprod{\phi_{\sigma}}{\phi_{\tau}}|
=\|\phi_{\sigma}\|_2^2=\|\phi\|_2^2$,}

\noindent sementara

$$\eqalign{\int_{I_{\tau}}w_{\sigma}
&=\mu J_{\sigma}\int_{x_{\sigma}-\bover12\mu I_{\sigma}}
  ^{x_{\sigma}+\bover12\mu I_{\sigma}}w((x-x_{\sigma})\mu J_{\sigma})dx
=\int_{-1/2}^{1/2}w,\cr}$$

\noindent sehingga tentu $|\innerprod{\phi_{\sigma}}{\phi_{\tau}}| \le
C_3\int_{I_{\tau}}w_{\sigma}$.

\medskip

\quad{\bf (ii) } Jika $\sigma\ne\tau$ dan $I_{\sigma}=I_{\tau}$ maka
$J_{\sigma}\cap J_{\tau}=\emptyset$, sehingga
$\innerprod{\phi_{\sigma}}{\phi_{\tau}}=0$, dengan 286E(b-iii).

\medskip

\quad{\bf (iii)} Sekarang anggap $I_{\sigma}\ne I_{\tau}$. Karena
$\mu I_{\tau}\le\mu I_{\sigma}$, seluruh $I_{\tau}$ harus terletak pada sisi
yang sama terhadap $x_{\sigma}$, sehingga $\int_{I_{\tau}}w_{\sigma}\ge
w_{\sigma}(x_{\tau})\mu I_{\tau}$, dengan (c). Dengan demikian

$$\eqalignno{|\innerprod{\phi_{\sigma}}{\phi_{\tau}}|
&\le\int_{-\infty}^{\infty}|\phi_{\sigma}|\times|\phi_{\tau}|
\le C_1^2\sqrt{\mu I_{\sigma}}\sqrt{\mu I_{\tau}}
  \int_{-\infty}^{\infty}w_{\sigma}\times w_{\tau}
\Displaycause{dengan menggunakan (e) dua kali}
=C_1^2\sqrt{\mu J_{\sigma}}\sqrt{\mu J_{\tau}}\int_{-\infty}^{\infty}
  w((x-x_{\sigma})\mu J_{\sigma})w((x-x_{\tau})\mu J_{\tau})dx\cr&
=C_1^2\sqrt{\mu J_{\sigma}}\sqrt{\mu I_{\tau}}\int_{-\infty}^{\infty}
  w(x\mu J_{\sigma}\mu I_{\tau}
    +(x_{\tau}-x_{\sigma})\mu J_{\sigma})w(x)dx\cr&
\le C_1^2C_2\sqrt{\mu J_{\sigma}}\sqrt{\mu I_{\tau}}
  w((x_{\tau}-x_{\sigma})\mu J_{\sigma})
\Displaycause{menurut (f), karena $\mu J_{\sigma}\mu I_{\tau}\le 1$}
\le C_3\sqrt{\mu I_{\sigma}}\sqrt{\mu I_{\tau}}w_{\sigma}(x_{\tau})
\le C_3\sqrt{\mu I_{\sigma}}\sqrt{\mu J_{\tau}}
   \int_{I_{\tau}}w_{\sigma},\cr}$$

\noindent sebagaimana diperlukan.

\medskip

{\bf (h)} Tetapkan $C_4=2\sum_{j=0}^{\infty}\int_{j+\bover12}^{\infty}w$; ini
adalah hingga karena $\int_{\alpha}^{\infty}w=\Bover1{2(1+\alpha)^2}$ untuk
setiap $\alpha\ge 0$.

Jika $k<k_{\tau}$, maka $k_{\sigma}\ne k$ untuk setiap $\sigma\ge\tau$,
sehingga hasilnya langsung berlaku. Jika $k\ge k_{\tau}$, maka untuk setiap
subinterval diadik $I$ dari $I_{\tau}$ yang panjangnya $2^{-k}$ terdapat
tepat satu $\sigma\ge\tau$ sedemikian sehingga $I_{\sigma}=I$, karena
$J_{\sigma}$ harus merupakan satu-satunya interval diadik sepanjang $2^k$
yang memuat $J_{\tau}$. Daftarkan semuanya sebagai
$\sigma_0,\ldots$ dalam urutan naik dari pusat-pusat $x_{\sigma_j}$, sehingga
jika $I_{\tau}=\coint{m\mu I_{\tau},(m+1)\mu I_{\tau}}$, maka
$x_{\sigma_j}=m\mu I_{\tau}+2^{-k}(j+\bover12)$, untuk
$j<2^{k-k_{\tau}}$.

$$\eqalign{\sum_{j=0}^{2^{k-k_{\tau}}-1}
  \int_{-\infty}^{m\mu I_{\tau}}w_{\sigma_j}
&=2^k\sum_{j=0}^{2^{k-k_{\tau}}-1}\int_{-\infty}^{m\mu I_{\tau}}
  w(2^k(x-m\mu I_{\tau})-j-\Bover12)dx\cr
&=\sum_{j=0}^{2^{k-k_{\tau}}-1}
  \int_{-\infty}^0w(x-j-\Bover12)dx\cr
&\le\sum_{j=0}^{\infty}
  \int_{j+\bover12}^{\infty}w
=\Bover12C_4.\cr}$$

\noindent Demikian pula (karena $w$ merupakan fungsi genap, sehingga seluruh
konfigurasi simetris terhadap $x_{\tau}$)

\Centerline{$\sum_{j=0}^{2^{k-k_{\tau}}-1} \int_{(m+1)\mu
I_{\tau}}^{\infty}w_{\sigma_j} \le\Bover12C_4$,}

\noindent dan

\Centerline{$\sum_{\sigma\ge\tau,k_{\sigma}=k} \int_{\Bbb R\setminus
I_{\tau}}w_{\sigma}\le C_4$,}

\noindent sebagaimana diperlukan. }%end of proof of 286G

\leader{286H}{`Massa' dan `energi'}\cmmnt{ ({\smc Lacey \& Thiele 00})}
Jika $P$
adalah himpunan bagian dari $Q$, $E\subseteq\Bbb R$ adalah terukur, $g:\Bbb
R\to\Bbb R$ terukur, dan $f\in\eusm L^2_{\Bbb C}$, tetapkan

\Centerline{$\mass_{Eg}(P) =\sup_{\sigma\in P,\tau\in Q,\tau\le\sigma}
\int_{E\cap g^{-1}[J_{\tau}]}w_{\tau} \le\sup_{\tau\in
Q}\int_{-\infty}^{\infty}w_{\tau}=1$,}

\Centerline{$\Delta_f(P) =\sum_{\sigma\in
P}|\innerprod{f}{\phi_{\sigma}}|^2$,}

\Centerline{$\energy_f(P) =\sup_{\tau\in Q}\sqrt{\mu
J_{\tau}}\sqrt{\Delta_f(P\cap T_{\tau})}$.}

\noindent Jika $P'\subseteq P$, maka
$\mass_{Eg}(P')\le\mass_{Eg}(P)$ dan
$\energy_f(P')\le\energy_f(P)$. \cmmnt{Perhatikan bahwa}
$\energy_f(\{\sigma\})=\sqrt{\mu J_{\sigma}}
|\innerprod{f}{\phi_{\sigma}}|$ untuk setiap $\sigma\in Q$\cmmnt{, karena
jika $\sigma\in T_{\tau}$, maka $\mu J_{\tau}\le\mu J_{\sigma}$}.
%end of 286H

\leader{286I}{Lema} Jika $P\subseteq Q$ berhingga dan $f\in\eusm
L^2_{\Bbb C}$, maka

(a) $\Delta_f(P)\le\|\sum_{\sigma\in P}
\innerprod{f}{\phi_{\sigma}}\phi_{\sigma}\|_2\|f\|_2$,

(b) $\sum_{\sigma,\tau\in P,J_{\sigma}=J_{\tau}}
\bigl|\innerprod{f}{\phi_{\sigma}} \innerprod{\phi_{\sigma}}{\phi_{\tau}}
\innerprod{\phi_{\tau}}{f}\bigr|\le C_3\Delta_f(P)$.

\proof{{\bf (a)}

\Centerline{$\Delta_f(P) =\sum_{\sigma\in P}
\innerprod{f}{\phi_{\sigma}}\innerprod{\phi_{\sigma}}{f}
=\Innerprod{\sum_{\sigma\in P} \innerprod{f}{\phi_{\sigma}}\phi_{\sigma}}{f}
\le\|\sum_{\sigma\in P} \innerprod{f}{\phi_{\sigma}}\phi_{\sigma}\|_2\|f\|_2$}

\noindent menurut ketaksamaan Cauchy (244Eb).

\medskip

{\bf (b)}

$$\eqalignno{\sum_{\Atop{\sigma,\tau\in P}{J_{\sigma}=J_{\tau}}}
  \bigl|\innerprod{f}{\phi_{\sigma}}
    \innerprod{\phi_{\sigma}}{\phi_{\tau}}
    \innerprod{\phi_{\tau}}{f}\bigr|
&\le\sum_{\Atop{\sigma,\tau\in P}{J_{\sigma}=J_{\tau}}}
  \Bover12\bigl(|\innerprod{f}{\phi_{\sigma}}|^2
             +|\innerprod{f}{\phi_{\tau}}|^2\bigr)
  |\innerprod{\phi_{\sigma}}{\phi_{\tau}}|\cr
\displaycause{karena $|\xi\zeta|\le\bover12(|\xi|^2+|\zeta|^2)$ untuk
semua bilangan kompleks $\xi$, $\zeta$}
&=\sum_{\sigma\in P}\sum_{\Atop{\tau\in P}{J_{\sigma}=J_{\tau}}}
  |\innerprod{f}{\phi_{\sigma}}|^2
  |\innerprod{\phi_{\sigma}}{\phi_{\tau}}|\cr
&\le\sum_{\sigma\in P}|\innerprod{f}{\phi_{\sigma}}|^2
  \sum_{\Atop{\tau\in P}{J_{\sigma}=J_{\tau}}}
     C_3\int_{I_{\tau}}w_{\sigma}\cr
\displaycause{menurut 286Gg, karena $k_{\sigma}=k_{\tau}$ jika
$J_{\sigma}=J_{\tau}$}
&\le\sum_{\sigma\in P}|\innerprod{f}{\phi_{\sigma}}|^2
  C_3\int_{-\infty}^{\infty}w_{\sigma}\cr
\displaycause{karena jika $\tau$, $\tau'$ anggota berbeda dari $P$ dan
$J_{\tau}=J_{\tau'}$, maka $I_{\tau}$ dan $I_{\tau'}$ saling lepas}
&=C_3\sum_{\sigma\in P}|\innerprod{f}{\phi_{\sigma}}|^2
=C_3\Delta_f(P).\cr}$$
}%end of proof of 286I

\leader{286J}{Lema} Tetapkan $C_5=2^{12}$. Jika $P\subseteq Q$ berhingga,
$E\subseteq\Bbb R$ terukur, $g:\Bbb R\to\Bbb R$ terukur, dan
$\gamma\ge\mass_{Eg}(P)$, maka kita bisa menemukan $R\subseteq Q$ sedemikian
sehingga $\gamma\sum_{\tau\in R}\mu I_{\tau}\le C_5\mu E$ dan\cmmnt{ (dalam
notasi dari 286Fb)} $\mass_{Eg}(P\setminus R^+)\le\bover14\gamma$.

\proof{{\bf (a)} Jika $\gamma=0$, kita dapat mengambil $R=\emptyset$. Jika
tidak, tetapkan $P_1=\{\sigma:\sigma\in P$,
$\mass_{Eg}(\{\sigma\})>\bover14\gamma\}$. Untuk setiap $\sigma\in P_1$, pilih
$\sigma'\in Q$ sehingga $\sigma'\le\sigma$ dan $\int_{E\cap
g^{-1}[J_{\sigma'}]}w_{\sigma'}>\bover14\gamma$. Misalkan $R$ ialah himpunan
anggota minimal $\{\sigma':\sigma\in P_1\}$ terhadap $\le$. Maka
$P\setminus R^+\subseteq\{\sigma:\mass_{Eg}(\{\sigma\})\le\bover14\gamma\}$
jadi $\mass_{Eg}(P\setminus R^+)\le\bover14\gamma$.

\medskip

{\bf (b)} Untuk $k\in\Bbb N$, tetapkan

\Centerline{$R_k =\{\tau:\tau\in R$, $\mu J_{\tau}\cdot\mu(E\cap
g^{-1}[J_{\tau}]\cap I^{(k)}_{\tau}) \ge 2^{2k-9}\gamma\}$,}

\noindent di mana $I^{(k)}_{\tau}$ adalah setengah terbuka interval dengan
pusat yang sama dengan $I_{\tau}$ dan $2^k$ dikalikan panjangnya. Sekarang
$R=\bigcup_{k\in\Bbb N}R_k$. \Prf\ Mengambil $\tau\in R$. Jika $k\in\Bbb N$
dan $x\in\Bbb R\setminus I^{(k)}_{\tau}$, maka $|x-x_{\tau}|\ge\bover12\mu
I^{(k)}_{\tau}=2^{k-1}\mu I_{\tau}$, jadi

\Centerline{$w_{\tau}(x) =w((x-x_{\tau})\mu J_{\tau})\mu J_{\tau} \le
w(2^{k-1})\mu J_{\tau} =(1+2^{k-1})^{-3}\mu J_{\tau}$.}

\noindent Dengan demikian

$$\eqalign{\Bover14\gamma
&<\int_{E\cap g^{-1}[J_{\tau}]}w_{\tau}
=\int_{E\cap g^{-1}[J_{\tau}]\cap I_{\tau}}w_{\tau}
  +\sum_{k=0}^{\infty}\int_{E\cap g^{-1}[J_{\tau}]
    \cap I^{(k+1)}_{\tau}\setminus I^{(k)}_{\tau}}w_{\tau}\cr
&\le \mu J_{\tau}\cdot\mu(E\cap g^{-1}[J_{\tau}]\cap I_{\tau})
  +\sum_{k=0}^{\infty}(1+2^{k-1})^{-3}\mu J_{\tau}\cdot
   \mu(E\cap g^{-1}[J_{\tau}]\cap I^{(k+1)}_{\tau}).\cr}$$

\noindent Ini berarti bahwa salah satu dari

\Centerline{$\mu J_{\tau}\cdot\mu(E\cap g^{-1}[J_{\tau}]\cap I_{\tau})
\ge\Bover18\gamma$}

\noindent dan $\tau\in R_0$, atau ada beberapa $k\in\Bbb N$ sedemikian
sehingga

\Centerline{$(1+2^{k-1})^{-3}\mu J_{\tau} \cdot\mu(E\cap g^{-1}[J_{\tau}]\cap
I^{(k+1)}_{\tau}) \ge 2^{-k-4}\gamma$}

\noindent dan

\Centerline{$\mu J_{\tau} \cdot\mu(E\cap g^{-1}[J_{\tau}]\cap
I^{(k+1)}_{\tau}) \ge(1+2^{k-1})^32^{-k-4}\gamma \ge 2^{2k-7}\gamma$,}

\noindent sehingga $\tau\in R_{k+1}$.\ \Qed

\medskip

{\bf (c)} Untuk setiap $k\in\Bbb N$, $\gamma\sum_{\tau\in R_k}\mu I_{\tau}\le
2^{11-k}\mu E$. \Prf\ Jika $R_k=\emptyset$, pernyataan ini langsung.
Jika tidak, daftarkan $R_k$ sebagai $\langle\tau_j\rangle_{j\le n}$ sedemikian rupa
sehingga $k_{\tau_j}\le k_{\tau_l}$ jika $j\le l\le n$. Definisikan
$q:\{0,\ldots,n\}\to\{0,\ldots,n\}$ secara induktif dengan aturan

\Centerline{$q(l) =\min(\{l\}\cup\{j:j<l$, $q(j)=j$, $(I^{(k)}_{\tau_j}\times
J_{\tau_j}) \cap(I^{(k)}_{\tau_l}\times J_{\tau_l})\ne\emptyset\})$}

\noindent untuk setiap $l\le n$. Perhatikan bahwa, untuk $l\le n$,
$q(q(l))=q(l)\le l$ dan $I^{(k)}_{\tau_{q(l)}}\cap
I^{(k)}_{\tau_l}\ne\emptyset$, sehingga bahwa

\Centerline{$I_{\tau_l}\subseteq I^{(k)}_{\tau_l} \subseteq
I^{(k+2)}_{\tau_{q(l)}}$, }

\noindent karena $\mu I_{\tau_l}^{(k)}\le\mu I_{\tau_{q(l)}}^{(k)}$. Selain
itu, jika $j<l\le n$ dan $q(j)=q(l)$, maka kedua $J_{\tau_j}$ dan $J_{\tau_l}$
bertemu $J_{\tau_{q(j)}}$, maka termasuklah itu, dan $J_{\tau_j}\subseteq
J_{\tau_l}$. Tapi karena $\tau_j$ dan $\tau_l$ adalah anggota yang berbeda
dari $R$, $\tau_j\not\le\tau_l$ dan $I_{\tau_j}\cap I_{\tau_l}$ harus kosong.

Tetapkan $M=\{q(j):j\le n\}$. Kita mempunyai

$$\eqalign{\gamma\sum_{\tau\in R_k}\mu I_{\tau}
&=\gamma\sum_{m\in M}\sum_{\Atop{j\le n}{q(j)=m}}\mu I_{\tau_j}
\le\gamma\sum_{m\in M}\mu I^{(k+2)}_{\tau_m}
=2^{k+2}\gamma\sum_{m\in M}\mu I_{\tau_m}\cr
&\le 2^{k+2}\sum_{m\in M}
  2^{9-2k}\mu(E\cap g^{-1}[J_{\tau_m}]\cap I^{(k)}_{\tau_m})\cr
&\le 2^{k+2}\cdot 2^{9-2k}\mu E
=2^{11-k}\mu E\cr}$$

\noindent karena jika $l$, $m\in M$ dan $l<m$ maka $I^{(k)}_{\tau_l}\times
J_{\tau_l}$ dan $I^{(k)}_{\tau_m}\times J_{\tau_m}$ saling lepas (karena
$q(m)=m$), sehingga $g^{-1}[J_{\tau_l}]\cap I^{(k)}_{\tau_l}$ dan
$g^{-1}[J_{\tau_m}]\cap I^{(k)}_{\tau_m}$ saling lepas.\ \Qed

\medskip

{\bf (d)} Oleh karena itu

\Centerline{$\gamma\sum_{\tau\in R}\mu I_{\tau}
\le\gamma\sum_{k=0}^{\infty}\sum_{\tau\in R_k}\mu I_{\tau} \le 2^{12}\mu E$,}

\noindent sebagaimana diperlukan. }%end of proof of 286J

\leader{286K}{Lema} Tetapkan $C_6=4(C_3+4C_3\sqrt{2C_4})$. Misalkan
$P\subseteq Q$ berhingga, $f\in\eusm L^2_{\Bbb C}$, $\|f\|_2=1$, dan
$\gamma\ge\energy_f(P)$. Maka kita dapat menemukan $R\subseteq Q$ sedemikian
sehingga
$\gamma^2\sum_{\tau\in R}\mu I_{\tau}\le C_6$ dan $\energy_f(P\setminus
R^+)\le\bover12\gamma$.

\proof{{\bf (a)} Kita boleh menganggap $\gamma>0$ dan $P\ne\emptyset$,
sebab jika tidak kita dapat mengambil $R=\emptyset$.

\medskip

\quad{\bf (i)} Hanya ada berhingga banyak himpunan berbentuk
$P\cap T_{\tau}$ untuk $\tau\in Q$; misalkan $\tilde R\subseteq Q$ suatu
himpunan berhingga tak kosong sedemikian sehingga, kapan pun $\tau\in Q$ dan
$P\cap T_{\tau}$ tak kosong, terdapat $\tau'\in \tilde R$ dengan
$P\cap T_{\tau}=P\cap T_{\tau'}$ dan $k_{\tau'}\ge
k_{\tau}$; ini mungkin karena jika $A\subseteq P$ tidak kosong maka
$k_{\tau}\le\min_{\sigma\in A}k_{\sigma}$ setiap kali $T_{\tau}\supseteq A$.

\medskip

\quad{\bf (ii) } Pilih $\tau_0,\tau_1,\ldots$ dan $P_0,P_1,\ldots$ secara
induktif, sebagai berikut. $P_0=P$. Karena $P_j\subseteq P$ tidak kosong,
pertimbangkan

\Centerline{$R_j =\{\tau:\tau\in\tilde R$, $\Delta_f(P_j\cap T_{\tau})
\ge\Bover14\gamma^2\mu I_{\tau}\}$.}

\noindent Jika $R_j=\emptyset$, hentikan induksi dan tetapkan $n=j$ serta
$R=\{\tau_l:l<j\}$. Jika tidak, pilih dari $R_j$ suatu anggota dengan
$y_{\tau}$ paling kiri, dan namai anggota itu $\tau_j$; tetapkan
$P_{j+1}=P_j\setminus\{\tau_j\}^+$, lalu lanjutkan. Perhatikan bahwa karena
$R_{j+1}\subseteq R_j$ untuk setiap $j$, $y_{\tau_{j+1}}\ge y_{\tau_j}$ untuk
setiap $j$.

Induksi harus berhenti setelah berhingga banyak tahap, sebab jika belum
berhenti pada $n=j$, maka $\Delta_f(P_j\cap T_{\tau_j})>0$, sehingga
$P_j\cap T_{\tau_j}$ tak kosong dan
$P_{j+1}\subseteq P_j\setminus T_{\tau_j}$ merupakan himpunan bagian sejati
dari $P_j$, sedangkan $P_0=P$ berhingga. Karena
$R_n=\emptyset$,

$$\eqalign{\energy_f(P\setminus R^+)
&=\energy_f(P_n)
=\sup_{\tau\in Q}\sqrt{\mu J_{\tau}}\sqrt{\Delta_f(P_n\cap T_{\tau})}\cr&
=\max_{\tau\in\tilde R}\sqrt{\mu J_{\tau}}\sqrt{\Delta_f(P_n\cap T_{\tau})}
\le\Bover12\gamma.\cr}$$

\medskip

\quad{\bf (iii)} Tetapkan
$P'_j=P_j\cap T_{\tau_j}\subseteq P_j\setminus P_{j+1}$ untuk $j<n$,
sehingga $\ofamily{j}{n}{P'_j}$ saling lepas, dan
$P'=\bigcup_{j<n}P'_j\subseteq P$. Kemudian jika $\sigma\in P'$, $j<n$ dan
$J_{\tau_j}\subseteq J^l_{\sigma}$, $I_{\sigma}\cap I_{\tau_j}=\emptyset$.
\Prf\ Misalkan $l<n$ memenuhi $\sigma\in P'_l$. $y_{\tau_j}\in
J_{\tau_j}\subseteq J^l_{\sigma}$ dan $y_{\tau_l}\in J^r_{\tau_l}\subseteq
J^r_{\sigma}$, sehingga $y_{\tau_j}<y_{\tau_l}$ dan $j<l$. Maka sesuai dengan
itu $P_{j+1}\supseteq P_l$ mengandung $\sigma$, sehingga
$\sigma\not\ge\tau_j$; karena $J_{\tau_j}\subseteq J_{\sigma}$,
$I_{\sigma}\not\subseteq I_{\tau_j}$, sedangkan $\mu I_{\sigma}\le\mu
I_{\tau_j}$, sehingga $I_{\sigma}$ saling lepas dengan $I_{\tau_j}$.\ \Qed

Oleh karena itu, jika $\sigma$, $\tau\in P'$ berbeda dan $J^l_{\sigma}\cap
J^l_{\tau}$ tidak kosong, maka $I_{\sigma}\cap I_{\tau}=\emptyset$. \Prf\ Jika
$J_{\sigma}=J_{\tau}$, hal ini benar karena $\sigma\ne\tau$. Jika tidak,
karena $J_{\sigma}$ dan $J_{\tau}$ beririsan, salah satunya termuat dalam
yang lain; anggap bahwa $J_{\sigma}\subset J_{\tau}$. Karena $J_{\sigma}$
bertemu $J^l_{\tau}$, $J_{\sigma}\subseteq J^l_{\tau}$. Sekarang misalkan
$j<n$ sehingga $\sigma\in P'_j$; maka $\sigma\ge\tau_j$, jadi
$J_{\tau_j}\subseteq J_{\sigma}\subseteq J^l_{\tau}$, dan $I_{\sigma}\cap
I_{\tau}\subseteq I_{\tau_j}\cap I_{\tau}=\emptyset$ menurut pengamatan terakhir.\
\Qed

\medskip

{\bf (b)} Sekarang mari kita memperkirakan

\Centerline{$\gamma^2\sum_{j<n}\mu I_{\tau_j} \le
4\sum_{j<n}\Delta_f(P'_j)=4\Delta_f(P')=4\alpha$}

\noindent dengan notasi seperti di atas. Karena $\|f\|_2=1$, kita
memiliki $\alpha\le\|\sum_{\sigma\in P'}
\innerprod{f}{\phi_{\sigma}}\phi_{\sigma}\|_2$ (286Ia). Jadi

$$\eqalignno{\alpha^2
&\le\|\sum_{\sigma\in P'}
  \innerprod{f}{\phi_{\sigma}}\phi_{\sigma}\|^2_2
=\sum_{\sigma,\tau\in P'}
  \innerprod{f}{\phi_{\sigma}}\innerprod{\phi_{\sigma}}{\phi_{\tau}}
  \innerprod{\phi_{\tau}}{f}\cr
&=\sum_{\Atop{\sigma,\tau\in P'}{J_{\sigma}=J_{\tau}}}
  \innerprod{f}{\phi_{\sigma}}\innerprod{\phi_{\sigma}}{\phi_{\tau}}
  \innerprod{\phi_{\tau}}{f}\cr
&\mskip100mu
 +\sum_{\Atop{\sigma,\tau\in P'}{J_{\sigma}\subseteq J^l_{\tau}}}
  \innerprod{f}{\phi_{\sigma}}\innerprod{\phi_{\sigma}}{\phi_{\tau}}
  \innerprod{\phi_{\tau}}{f}
 +\sum_{\Atop{\sigma,\tau\in P'}{J_{\tau}\subseteq J^l_{\sigma}}}
  \innerprod{f}{\phi_{\sigma}}\innerprod{\phi_{\sigma}}{\phi_{\tau}}
  \innerprod{\phi_{\tau}}{f}\cr}$$

\noindent karena $\innerprod{\phi_{\sigma}}{\phi_{\tau}}=0$ kecuali
$J^l_{\sigma}\cap J^l_{\tau}\ne\emptyset$, seperti yang dinyatakan dalam
286E(b-iii).

Tinjau ketiga suku ini secara terpisah. Untuk suku pertama, kita memiliki

\Centerline{$\sum_{\sigma,\tau\in P',J_{\sigma}=J_{\tau}}
\bigl|\innerprod{f}{\phi_{\sigma}} \innerprod{\phi_{\sigma}}{\phi_{\tau}}
\innerprod{\phi_{\tau}}{f}\bigr| \le C_3\alpha$}

\noindent menurut 286Ib. Untuk suku kedua, kita memiliki

$$\eqalign{\sum_{\Atop{\sigma,\tau\in P'}{J_{\sigma}\subseteq J^l_{\tau}}}
  \bigl|\innerprod{f}{\phi_{\sigma}}
   \innerprod{\phi_{\sigma}}{\phi_{\tau}}
   \innerprod{\phi_{\tau}}{f}\bigr|
&\le\sum_{\sigma\in P'}|\innerprod{f}{\phi_{\sigma}}|
  \sum_{\Atop{\tau\in P'}{J_{\sigma}\subseteq J^l_{\tau}}}
   |\innerprod{\phi_{\sigma}}{\phi_{\tau}}
   \innerprod{\phi_{\tau}}{f}|\cr
&\le\sqrt{\sum_{\sigma\in P'}|\innerprod{f}{\phi_{\sigma}}|^2}
  \sqrt{\sum_{\sigma\in P'}
   \bigl(\sum_{\Atop{\tau\in P'}{J_{\sigma}\subseteq J^l_{\tau}}}
     |\innerprod{\phi_{\sigma}}{\phi_{\tau}}
     \innerprod{\phi_{\tau}}{f}|\bigr)^2}\cr
&=\sqrt{\alpha}\sqrt{\sum_{j<n}H_j},\cr}$$

\noindent dengan, untuk $j<n$, kita tetapkan

$$H_j=\sum_{\sigma\in P'_j}
   \bigl(\sum_{\Atop{\tau\in P'}{J_{\sigma}\subseteq J^l_{\tau}}}
     |\innerprod{\phi_{\sigma}}{\phi_{\tau}}
     \innerprod{\phi_{\tau}}{f}|\bigr)^2.$$

\noindent Sekarang kita bisa memperkirakan $H_j$ dengan mengamati bahwa, untuk
setiap $\tau\in P'$,

\Centerline{$|\innerprod{\phi_{\tau}}{f}| =\sqrt{\mu
I_{\tau}}\energy_f(\{\tau\})\le\gamma\sqrt{\mu I_{\tau}}$,}

\noindent sedangkan jika $\sigma$, $\tau\in P'$ dan $J_{\tau}^l\supseteq
J_{\sigma}$ maka

\Centerline{$|\innerprod{\phi_{\sigma}}{\phi_{\tau}}| \le C_3\sqrt{\mu
I_{\sigma}}\sqrt{\mu J_{\tau}}\int_{I_{\tau}}w_{\sigma}$}

\noindent menurut 286Gg. Kita juga perlu mengetahui bahwa jika $\sigma\in P'_j$ dan
$\tau$, $\tau'$ adalah anggota berbeda dari $P'$ sedemikian sehingga
$J_{\sigma}\subseteq J^l_{\tau}\cap J^l_{\tau'}$, maka $I_{\tau}$, $I_{\tau'}$
dan $I_{\tau_j}$ saling lepas, menurut (a-iii) di atas, karena
$J_{\tau_j}\subseteq J_{\sigma}$. Jadi kita memiliki

$$\eqalignno{H_j
&\le\sum_{\sigma\in P'_j}
   \bigl(\sum_{\Atop{\tau\in P'}{J_{\sigma}\subseteq J^l_{\tau}}}
    \gamma\sqrt{\mu I_{\tau}}
      \cdot C_3\sqrt{\mu I_{\sigma}}\sqrt{\mu J_{\tau}}
      \int_{I_{\tau}}w_{\sigma}\bigr)^2\cr
&=C_3^2\gamma^2\sum_{\sigma\in P'_j}\mu I_{\sigma}
   \bigl(\sum_{\Atop{\tau\in P'}{J_{\sigma}\subseteq J^l_{\tau}}}
     \int_{I_{\tau}}w_{\sigma}\bigr)^2
\le C_3^2\gamma^2\sum_{\sigma\in P'_j}\mu I_{\sigma}
   (\int_{\Bbb R\setminus I_{\tau_j}}w_{\sigma})^2\cr
&\le C_3^2\gamma^2\sum_{k=k_{\tau_j}}^{\infty}2^{-k}
   \sum_{\Atop{\sigma\in P'_j}{k_{\sigma}=k}}
   \int_{\Bbb R\setminus I_{\tau_j}}w_{\sigma}
   \cdot\int_{-\infty}^{\infty}w_{\sigma}
\le C_3^2\gamma^2\sum_{k=k_{\tau_j}}^{\infty}2^{-k}C_4\cr
\displaycause{menurut 286Ga dan 286Gh, karena $\sigma\ge\tau_j$ untuk setiap
$\sigma\in P'_j$}
&=C_3^2\gamma^22^{-k_{\tau_j}+1}C_4
=2C_3^2C_4\gamma^2\mu I_{\tau_j}.\cr}$$

\noindent Dengan demikian

\Centerline{$\sum_{j<n}H_j \le 2C_3^2\gamma^2C_4\sum_{j<n}\mu I_{\tau_j} \le
2C_3^2C_4\cdot 4\alpha$,}

\noindent dan

\Centerline{$\sum_{\sigma,\tau\in P',J_{\sigma}\subseteq J^l_{\tau}}
|\innerprod{f}{\phi_{\sigma}} \innerprod{\phi_{\sigma}}{\phi_{\tau}}
\innerprod{\phi_{\tau}}{f}| \le\sqrt{\alpha\sumop_{j<n}H_j} \le
2C_3\alpha\sqrt{2C_4}$.}

Demikian juga,

\Centerline{$\sum_{\sigma,\tau\in P',J_{\tau}\subseteq J^l_{\sigma}}
|\innerprod{f}{\phi_{\sigma}} \innerprod{\phi_{\sigma}}{\phi_{\tau}}
\innerprod{\phi_{\tau}}{f}| \le 2C_3\alpha\sqrt{2C_4}$;}

\noindent dengan menggabungkan hasil-hasil ini,

\Centerline{$\alpha^2 \le\alpha(C_3+4C_3\sqrt{2C_4}) =\Bover14\alpha C_6$}

\noindent dan $\alpha\le\bover14C_6$. Ini berarti bahwa

\Centerline{$\gamma^2\sum_{j<n}\mu I_{\tau_j} \le 4\alpha\le C_6$,}

\noindent dan $R=\{\tau_j:j<n\}$ memiliki kedua sifat yang diperlukan. }%end of proof of 286K

\leader{286L}{Lema} Tetapkan

\Centerline{$C_7 =C_1\bigl(\Bover{7}2+\Bover87+\Bover{28}{w(3/2)}
+\Bover{4\sqrt{14 C_3}}{w(3/2)}\bigr)$.}

\noindent Misalkan $P$ merupakan himpunan bagian berhingga dari $Q$ yang
memiliki batas bawah $\tau$ dalam $Q$ terhadap urutan $\le$, misalkan
$E\subseteq\Bbb R$ terukur, $g:\Bbb R\to\Bbb R$ terukur, dan
$f\in\eusm L^2_{\Bbb C}$. Maka

\Centerline{$\sum_{\sigma\in P}|\innerprod{f}{\phi_{\sigma}} \int_{E\cap
g^{-1}[J^r_{\sigma}]}\phi_{\sigma}| \le C_7\energy_f(P)\mass_{Eg}(P)\mu
I_{\tau}$.}

\proof{ Tetapkan $\gamma=\energy_f(P)$ dan $\gamma'=\mass_{Eg}(P)$. Jika
$P=\emptyset$, hasilnya langsung berlaku; jadi anggap $P\ne\emptyset$.

\medskip

{\bf (a)(i)} Perhatikan bahwa $\bigcup_{\sigma\in P}I_{\sigma}\subseteq
I_{\tau}$, $J_{\tau}\subseteq\bigcap_{\sigma\in P}J_{\sigma}$ dan
$k_{\tau}\le\min_{\sigma\in P}k_{\sigma}$. Jadi jika $\sigma$, $\sigma'\in P$
berbeda dan $\mu I_{\sigma}=\mu I_{\sigma'}$, maka $J_{\sigma}=J_{\sigma'}$
dan $I_{\sigma}\cap I_{\sigma'}=\emptyset$.

\medskip

\quad{\bf (ii) } Untuk interval diadik $I$, misalkan $I^*$ ialah interval
setengah terbuka yang berpusat sama dengan $I$ dan panjangnya tiga kali
panjang interval tersebut. Misalkan $\Cal J$ ialah keluarga semua $I\in\Cal I$ sedemikian
sehingga $I_{\sigma}\not\subseteq I^*$ untuk setiap $\sigma\in P$ dengan
$\mu I_{\sigma}\le\mu I$. Karena $P$ berhingga, semua interval yang cukup
kecil termasuk dalam $\Cal J$, dan $\bigcup\Cal J=\Bbb R$; misalkan $\Cal K$
ialah himpunan anggota maksimal $\Cal J$, sehingga anggota-anggota $\Cal K$
saling lepas. Kita mempunyai $\bigcup\Cal K=\Bbb R$. \Prf\ Karena
$P\ne\emptyset$, pilih sementara $\sigma\in P$. Untuk $I\in\Cal J$, bagi
setiap $n\in\Bbb N$ tinjau interval $\tilde I^{(n)}\in\Cal I$ yang memuat
$I$ dan panjangnya $2^n\mu I$. Ada $n\in\Bbb N$ sedemikian sehingga
$\mu\tilde I^{(n)}\ge\mu I_{\sigma}$ dan $I_{\sigma}\subseteq(\tilde
I^{(n)})^*$, sehingga $\tilde I^{(k)}\notin\Cal J$ untuk setiap $k\ge n$;
karena itu ada $k<n$ sedemikian sehingga $\tilde I^{(k)}\in\Cal K$. Jadi
$I\subseteq\tilde I^{(k)}\subseteq\bigcup\Cal K$; karena $I$ sembarang,
$\bigcup\Cal K=\bigcup\Cal J=\Bbb R$.\ \Qed

\medskip

\quad{\bf (iii)} Untuk $K\in\Cal K$, pilih $l_K\in\Bbb Z$ sehingga
$\mu K=2^{-l_K}$. Jika $l_K\ge k_{\tau}$, yaitu $\mu K\le\mu I_{\tau}$, maka
$K$ harus termuat dalam interval setengah terbuka $\hat I$ yang berpusat di
$x_{\tau}$ dan panjang $7\mu I_{\tau}$, karena jika tidak kita harus memiliki
$I_{\tau}\cap\tilde K^*=\emptyset$, dengan $\tilde K$ interval diadik
berpanjang $2\mu K$ yang memuat $K$, dan $\tilde K$ akibatnya termasuk dalam
$\Cal J$. Jadi

\Centerline{$\sum_{K\in\Cal K,\mu K\le\mu I_{\tau}}\mu K \le\mu\hat I=7\mu
I_{\tau}$,}

\noindent karena anggota-anggota $\Cal K$ saling lepas.

\medskip

\quad{\bf (iv) } Untuk setiap $l<k_{\tau}$, hanya ada tiga anggota $K$ dari
$\Cal K$ dengan $l_K=l$. \Prf\ Jika $I\in\Cal I$ dan $\mu I>\mu I_{\tau}$,
maka baik $I_{\tau}\subseteq I^*$ atau $I_{\tau}\cap I^*=\emptyset$, dan
$I\in\Cal J$ jika dan hanya jika $I_{\tau}\cap I^*$ kosong. Ini berarti bahwa
jika $K\in\Cal I$ dan $\mu K=2^{-l}$, maka $K\in\Cal K$ jika dan hanya jika
$I_{\tau}\cap K^*$ kosong dan
$I_{\tau}\subseteq\tilde K^*$. Jadi jika
$I_{\tau}\subseteq\coint{2^{-l}n,2^{-l}(n+1)}$ dan
$K=\coint{2^{-l}m,2^{-l}(m+1)}$, kita mempunyai $K\in\Cal K$ jika dan hanya jika

\Centerline{{\it baik} $m=n-2$ {\it atau} $m=n+2$ {\it atau} $m=n-3$ adalah
genap {\it atau} $m=n+3$ adalah ganjil;}

\noindent yang, untuk setiap $n$ tertentu, terjadi hanya bagi tiga nilai $m$.\
\Qed

\medskip

{\bf (b)} Untuk $\sigma\in P$, pilih bilangan kompleks $\zeta_{\sigma}$
bermodulus $1$ sedemikian sehingga
$\zeta_{\sigma}\innerprod{f}{\phi_{\sigma}} \int_{E\cap
g^{-1}[J^r_{\sigma}]}\phi_{\sigma}$ bernilai riil dan tak negatif. Tetapkan
$W=P\times\Cal K$. Untuk $(\sigma,K)\in W$, tetapkan

\Centerline{$\alpha_{\sigma K} =\innerprod{f}{\phi_{\sigma}} \int_{E\cap
g^{-1}[J^r_{\sigma}]\cap K}\phi_{\sigma}$.}

\noindent Tujuan dari bukti adalah untuk memperkirakan

\Centerline{$\sum_{\sigma\in P}\bigl|\innerprod{f}{\phi_{\sigma}} \int_{E\cap
g^{-1}[J^r_{\sigma}]}\phi_{\sigma}\bigr| =\sum_{(\sigma,K)\in
W}\zeta_{\sigma}\alpha_{\sigma K}$.}

\noindent Akan membantu untuk segera mencatat bahwa

\Centerline{$\sum_{(\sigma,K)\in W}|\alpha_{\sigma K}| \le\sum_{\sigma\in
P}|\innerprod{f}{\phi_{\sigma}}| \int_{-\infty}^{\infty}|\phi_{\sigma}|$}

\noindent adalah hingga.

Tetapkan

\Centerline{$W_0 =\{(\sigma,K):\sigma\in P$, $K\in\Cal K$, $\mu
I_{\sigma}\le\mu K\le\mu I_{\tau}\}$,}

\Centerline{$W_1 =\{(\sigma,K):\sigma\in P$, $K\in\Cal K$, $\mu I_{\tau}<\mu
K\}$,}

\Centerline{$W_2 =\{(\sigma,K):\sigma\in P\setminus T_{\tau}$, $K\in\Cal K$,
$\mu K<\mu I_{\sigma}\}$,}

\Centerline{$W_3 =\{(\sigma,K):\sigma\in P\cap T_{\tau}$, $K\in\Cal K$, $\mu
K<\mu I_{\sigma}\}$.}

\noindent Karena $\mu I_{\sigma}\le\mu I_{\tau}$ untuk setiap $\sigma\in P$,
$W=W_0\cup W_1\cup W_2\cup W_3$. Kita akan menaksir

\Centerline{$\alpha_j =\sum_{(\sigma,K)\in W_j}\zeta_{\sigma}\alpha_{\sigma
K}$}

\noindent untuk setiap $j$; empat komponen dalam ekspresi untuk $C_7$ yang
diberikan di atas adalah batas untuk $|\alpha_0|, |\alpha_1|$, $|\alpha_2|$
dan $|\alpha_3|$, berturut-turut.

\wheader{286L}{6}{2}{2}{36pt}

{\bf (c)(i)} Jika $K\in\Cal K$ dan $l_K=l$, maka untuk setiap $k\ge l$

\Centerline{$\sum_{\sigma\in P,k_{\sigma}=k}|\alpha_{\sigma K}| \le
2^{-k}C_1\gamma\gamma'(1+2^{k-l})^{-2} \le 2^{-k-2}C_1\gamma\gamma'$.}

\noindent\Prf\ Untuk setiap $\sigma\in P$,

\Centerline{$|\innerprod{f}{\phi_{\sigma}}| =\sqrt{\mu
I_{\sigma}}\energy_f(\{\sigma\}) \le\gamma\sqrt{\mu I_{\sigma}}$}

\noindent seperti yang dicatat dalam 286H, dan

$$\eqalignno{\int_{E\cap g^{-1}[J^r_{\sigma}]\cap K}|\phi_{\sigma}|
&\le C_1\mu I_{\sigma}\sqrt{\mu I_{\sigma}}
  \int_{E\cap g^{-1}[J^r_{\sigma}]\cap K}w_{\sigma}^2\cr
\displaycause{286Ge}
&\le C_1\mu I_{\sigma}\sqrt{\mu I_{\sigma}}
  \int_{E\cap g^{-1}[J_{\sigma}]}w_{\sigma}
  \cdot\sup_{x\in K}w_{\sigma}(x)\cr
&\le C_1\mu I_{\sigma}\sqrt{\mu I_{\sigma}}\gamma'
  \sup_{x\in K}w_{\sigma}(x)
=C_1\sqrt{\mu I_{\sigma}}\gamma'w(\mu J_{\sigma}\rho(x_{\sigma},K)),\cr}$$

\noindent dengan $\rho(x_{\sigma},K)$ menyatakan $\inf_{x\in
K}|x-x_{\sigma}|$. Jadi, untuk $k\ge l$,

$$\eqalign{\sum_{\Atop{\sigma\in P}{k_{\sigma}=k}}
  |\alpha_{\sigma K}|
&\le\sum_{\Atop{\sigma\in P}{k_{\sigma}=k}}
  C_1\gamma\gamma'\mu I_{\sigma}w(\mu J_{\sigma}\rho(x_{\sigma},K))\cr
&=2^{-k}C_1\gamma\gamma'
  \sum_{\Atop{\sigma\in P}{k_{\sigma}=k}}
    w(2^k\rho(x_{\sigma},K))
\le 2^{-k}C_1\gamma\gamma'
  \cdot 2\sum_{n=2^{k-l}}^{\infty}w(n+\hbox{$\bover12$})\cr}$$

\noindent karena $x_{\sigma}$, untuk $\sigma\in P$ dan $k_{\sigma}=k$, semua
berbeda (lihat (a-i) di atas) dan semua jarak setidaknya $\mu K=2^{k-l}2^{-k}$
dari $K$ (karena $I_{\sigma}\not\subseteq K^*$); jadi ada paling banyak dua
$\sigma$ sedemikian sehingga $\rho(x_{\sigma},K)=2^{-k}(n+\bover12)$
untuk setiap $n\ge 2^{k-l}$. Jadi kita memiliki

\Centerline{$\sum_{\sigma\in P,k_{\sigma}=k} |\alpha_{\sigma K}| \le
2^{-k}C_1\gamma\gamma'(1+2^{k-l})^{-2} \le 2^{-k-2}C_1\gamma\gamma'$}

\noindent oleh 286Gb.\ \Qed

\medskip

\quad{\bf (ii)} Sekarang

$$\eqalign{|\alpha_0|
&\le\sum_{(\sigma,K)\in W_0}|\alpha_{\sigma K}|
=\sum_{\Atop{K\in\Cal K}{\mu K\le\mu I_{\tau}}}
  \sum_{\Atop{\sigma\in P}{\mu I_{\sigma}\le\mu K}}
   |\alpha_{\sigma K}|\cr
&=\sum_{\Atop{K\in\Cal K}{\mu K\le\mu I_{\tau}}}\sum_{k=l_K}^{\infty}
   \sum_{\Atop{\sigma\in P}{k_{\sigma}=k}}
   |\alpha_{\sigma K}|
\le\sum_{\Atop{K\in\Cal K}{\mu K\le\mu I_{\tau}}}\sum_{k=l_K}^{\infty}
   2^{-k-2}C_1\gamma\gamma'\cr
&=C_1\gamma\gamma'\sum_{\Atop{K\in\Cal K}{\mu K\le\mu I_{\tau}}}
  2^{-l_K-1}
=\Bover12C_1\gamma\gamma'
   \sum_{\Atop{K\in\Cal K}{\mu K\le\mu I_{\tau}}}\mu K
\le\Bover{7}2C_1\gamma\gamma'\mu I_{\tau}\cr}$$

\noindent menurut rumus dalam (a-iii). Ini menangani $\alpha_0$.

\medskip

{\bf (d)} Selanjutnya pertimbangkan $W_1$. Kita memiliki

$$\eqalignno{|\alpha_1|
&\le\sum_{(\sigma,K)\in W_1}|\alpha_{\sigma K}|
=\sum_{l=-\infty}^{k_{\tau}-1}\sum_{k=k_{\tau}}^{\infty}
   \sum_{\Atop{K\in\Cal K}{l_K=l}}\sum_{\Atop{\sigma\in P}{k_{\sigma}=k}}
   |\alpha_{\sigma K}|\cr
&\le\sum_{l=-\infty}^{k_{\tau}-1}\sum_{k=k_{\tau}}^{\infty}
   \sum_{\Atop{K\in\Cal K}{l_K=l}}
   2^{-k}C_1\gamma\gamma'(1+2^{k-l})^{-2}\cr
\displaycause{menurut (c-i) di atas}
&=3C_1\gamma\gamma'\sum_{l=-\infty}^{k_{\tau}-1}\sum_{k=k_{\tau}}^{\infty}
   2^{-k}(1+2^{k-l})^{-2}\cr
\displaycause{menurut (a-iv)}
&\le 3C_1\gamma\gamma'\sum_{l=-\infty}^{k_{\tau}-1}
  \sum_{k=k_{\tau}}^{\infty}2^{-3k}2^{2l}
=3C_1\gamma\gamma'2^{2(k_{\tau}-1)}2^{-3k_{\tau}}
  \sum_{l=0}^{\infty}2^{-2l}\sum_{k=0}^{\infty}2^{-3k}\cr
&=\Bover34C_1\gamma\gamma'2^{-k_{\tau}}\cdot\Bover43\cdot\Bover87
=\Bover87C_1\gamma\gamma'\mu I_{\tau}.\cr}$$

\noindent Ini menangani $\alpha_1$.

\medskip

{\bf (e)} Untuk $K\in\Cal K$, tetapkan $G_K =K\cap E\cap\bigcup_{\sigma\in P,\mu
I_{\sigma}>\mu K}g^{-1}[J_{\sigma}]$. Maka $\mu G_K\le 2\gamma'\mu
K/w(\bover32)$. \Prf\ Jika $\mu I_{\tau}\le\mu K$, maka $G_K=\emptyset$, jadi
kita dapat menganggap $\mu K<\mu I_{\tau}$. Misalkan $\tilde K\in\Cal I$
ialah interval diadik yang memuat $K$ dan berpanjang dua kali lipat.
$\tilde K\notin\Cal J$, sehingga ada $\sigma\in P$ sedemikian sehingga
$\tilde K^*\supseteq I_{\sigma}$ dan

\Centerline{$\mu I_{\sigma}\le\mu\tilde K=2\mu K\le\mu I_{\tau}$.}

\noindent Pilih $\upsilon\in Q$ sehingga $\tau\le\upsilon\le\sigma$
dan $\mu I_{\upsilon}=2\mu K$ (286F(a-iv)). Maka $I_{\upsilon}$ bertemu
$\tilde K^*$, sehingga $\tilde K$ sama dengan $I_{\upsilon}$ atau berdekatan
dengannya, dan $|x-x_{\upsilon}|\le\bover32\cdot\mu I_{\upsilon}$ untuk setiap
$x\in\tilde K$, oleh karena itu untuk setiap $x\in K$.

\Centerline{$w_{\upsilon}(x)\ge w(\bover32)\mu J_{\upsilon} =w(\bover32)/2\mu
K$}

\noindent untuk setiap $x\in K$. Di sisi lain, karena $\sigma\in P$ dan
$\upsilon\le\sigma$, $\int_{E\cap
g^{-1}[J_{\upsilon}]}w_{\upsilon}\le\gamma'$. Jadi

\Centerline{$\mu(E\cap g^{-1}[J_{\upsilon}]\cap K) \le 2\gamma'\mu
K/w(\bover32)$.}

Sekarang misalkan $\sigma'\in P$ dan $\mu I_{\sigma'}>\mu K$. Maka
$k_{\sigma'}\le k_{\upsilon}$ dan $J_{\sigma'}$ adalah interval diadik
berpanjang $2^{k_{\sigma'}}$ yang memuat $J_{\tau}$. Tetapi $J_{\upsilon}$ adalah
interval diadik berpanjang $2^{k_{\upsilon}}$ yang memuat $J_{\tau}$, sehingga
termuat dalam $J_{\sigma'}$, dan $g^{-1}[J_{\sigma'}]\subseteq
g^{-1}[J_{\upsilon}]$. Karena $\sigma'$ sembarang, $G_K\subseteq
E\cap g^{-1}[J_{\upsilon}]\cap K$ dan $\mu G_K\le 2\gamma'\mu K/w(\bover32)$,
seperti yang diklaim.\ \Qed

\medskip

{\bf (f)(i)} Jika $\sigma$, $\upsilon\in P\setminus T_{\tau}$ dan
$k_{\sigma}\ne k_{\upsilon}$, maka $J^r_{\sigma}\cap
J^r_{\upsilon}=\emptyset$. \Prf\ Cukup untuk mempertimbangkan kasus $\mu
J_{\sigma}<\mu J_{\upsilon}$, sehingga $\mu J_{\sigma}\le\mu J^r_{\upsilon}$.
Karena $J_{\sigma}$ termuat dalam $J_{\tau}$, sedangkan $J^r_{\upsilon}$ tidak,
$J_{\sigma}$ saling lepas dengan $J^r_{\upsilon}$, dan hasilnya mengikuti.\ \Qed

\medskip

\quad{\bf (ii)} Untuk $x\in\Bbb R$, tetapkan

\Centerline{$v_2(x)=\bigl|\sum_{(\sigma,K)\in W_2}
\zeta_{\sigma}\innerprod{f}{\phi_{\sigma}}\phi_{\sigma}(x) \chi(E\cap
g^{-1}[J^r_{\sigma}]\cap K)(x)\bigr|$.}

\noindent (Jumlah ini berhingga karena ada paling banyak satu $K\in\Cal K$
yang mengandung $x$.) Kemudian untuk setiap $x\in\Bbb R$ ada sebuah $k\ge
k_{\tau}$ sedemikian sehingga

\Centerline{$v_2(x) =\bigl|\sum_{\sigma\in P,k_{\sigma}=k}
\zeta_{\sigma}\innerprod{f}{\phi_{\sigma}}\phi_{\sigma}(x)\bigl|$.}

\noindent\Prf\ Jika $v_2(x)=0$, ambil sembarang $k$ yang cukup besar. Jika
tidak, $x\in E$ dan
kita punya pasangan $(\upsilon,L)\in W_2$ sedemikian sehingga $x\in
g^{-1}[J^r_{\upsilon}]\cap L$. Ambil $k=k_{\upsilon}$. Karena $L$ adalah
satu-satunya anggota $\Cal K$ yang mengandung $x$, jadi

\Centerline{$v_2(x) =\bigl|\sum_{\sigma\in P_x}
\innerprod{f}{\phi_{\sigma}}\phi_{\sigma}(x)\bigr|$,}

\noindent di mana $P_x=\{\sigma:\sigma\in P\setminus T_{\tau}$, $\mu
I_{\sigma}>\mu L$, $g(x)\in J^r_{\sigma}\}$. Sekarang jika $\sigma\in P$ dan
$k_{\sigma}=k$, maka $\mu I_{\sigma}=\mu I_{\upsilon}>\mu L$,
$J_{\sigma}=J_{\upsilon}$ dan $J^r_{\sigma}=J^r_{\upsilon}$ tidak termasuk
$J^r_{\tau}$, sehingga $\sigma\in P\setminus T_{\tau}$, $g(x)\in J^r_{\sigma}$
dan $\sigma\in P_x$. Di sisi lain, (i) di atas memberitahu kita bahwa
$k_{\sigma}=k$ setiap kali $\sigma\in P\setminus T_{\tau}$ dan $g(x)\in
J^r_{\sigma}$. Jadi $P_x=\{\sigma:\sigma\in P$, $k_{\sigma}=k\}$ dan

\Centerline{$v_2(x) =\bigl|\sum_{\sigma\in P,k_{\sigma}=k}
\zeta_{\sigma}\innerprod{f}{\phi_{\sigma}}\phi_{\sigma}(x)\bigl|$. \Qed}

\medskip

\quad{\bf (iii) } Akibatnya $v_2(x)\le 2C_1\gamma$ untuk setiap $x\in\Bbb R$.
\Prf\ Jika $v_2(x)=0$, hasilnya langsung. Jika tidak, ambil $k$ dari (ii).
Kemudian

$$\eqalignno{v_2(x)
&\le\sum_{\Atop{\sigma\in P}{k_{\sigma}=k}}
  |\innerprod{f}{\phi_{\sigma}}\phi_{\sigma}(x)|
\le\sum_{\Atop{\sigma\in P}{k_{\sigma}=k}}\sqrt{\mu I_{\sigma}}\gamma
  \cdot\sqrt{\mu I_{\sigma}}C_1w_{\sigma}(x)\cr
\displaycause{menurut 286H dan 286Ge}
&=C_1\gamma\sum_{\Atop{\sigma\in P}{k_{\sigma}=k}}w(2^k(x-x_{\sigma}))
\le C_1\gamma\sum_{n=-\infty}^{\infty}w(2^kx-n-\hbox{$\bover12$})\cr
\displaycause{karena semua $x_{\sigma}$, untuk $\sigma\in P$ dan
$k_{\sigma}=k$, berbeda dan berbentuk $2^{-k}(n+\bover12)$}
&\le 2C_1\gamma\cr}$$

\noindent oleh 286Gd.\ \Qed

\medskip

\quad{\bf (iv)} Perhatikan juga bahwa, jika $v_2(x)>0$, ada sepasang
$(\sigma,K)\in W_2$ sedemikian sehingga $x\in g^{-1}[J_{\sigma}]\cap K$;
jadi $\mu K<\mu I_{\sigma}\le\mu I_{\tau}$ dan $x\in G_K$. Sekarang

$$\eqalignno{|\alpha_2|
&=\bigl|\sum_{(\sigma,K)\in W_2}
  \zeta_{\sigma}\innerprod{f}{\phi_{\sigma}}\int_{-\infty}^{\infty}
  \phi_{\sigma}\times\chi(E\cap g^{-1}[J^r_{\sigma}]\cap K)\bigr|\cr
&\le\int_{-\infty}^{\infty}v_2
\le\sum_{\Atop{K\in\Cal K}{\mu K<\mu I_{\tau}}}\int_{G_K}v_2
\le\sum_{\Atop{K\in\Cal K}{\mu K<\mu I_{\tau}}}
  \bover{4C_1\gamma\gamma'\mu K}{w(\bover32)}\cr
\displaycause{dengan menggabungkan taksiran dalam (e) dan (iii) di atas}
&\le\bover{28\cdot C_1\gamma\gamma'\mu I_{\tau}}{w(\bover32)}\cr}$$

\noindent oleh (a-iii). Ini berkaitan dengan $\alpha_2$.

\medskip

{\bf (g)} Tetapkan $P'=P\cap T_{\tau}$ dan $\tilde f=\sum_{\sigma\in P'}
\zeta_{\sigma}\innerprod{f}{\phi_{\sigma}}\phi_{\sigma}$. Kemudian

\Centerline{$\|\tilde f\|^2_2 \le C_3\gamma^2\mu I_{\tau}$.}

\noindent\Prf\ Jika $\sigma$, $\sigma'\in P'$ dan $k_{\sigma}\ne
k_{\sigma'}$, maka $\innerprod{\phi_{\sigma}}{\phi_{\sigma'}}=0$ (286Fc).
Sementara jika $k_{\sigma}=k_{\sigma'}$, maka $J_{\sigma}=J_{\sigma'}$, oleh
(a-i). Jadi

$$\eqalignno{\|\tilde f\|_2^2
&=\sum_{\sigma,\sigma'\in P'}
  \zeta_{\sigma}\innerprod{f}{\phi_{\sigma}}
  \innerprod{\phi_{\sigma}}{\phi_{\sigma'}}
  \innerprod{\phi_{\sigma'}}{f}\bar\zeta_{\sigma'}\cr
&\le\sum_{\Atop{\sigma,\sigma'\in P'}{J_{\sigma}=J_{\sigma'}}}
  \bigl|\innerprod{f}{\phi_{\sigma}}
  \innerprod{\phi_{\sigma}}{\phi_{\sigma'}}
  \innerprod{\phi_{\sigma'}}{f}\bigr|
\le C_3\Delta_f(P')\cr
\displaycause{286Ib}
&\le C_3\gamma^2\mu I_{\tau}\cr}$$

\noindent menurut definisi `energi', karena $P'\subseteq T_{\tau}$.\ \Qed

\medskip

{\bf (h)} Untuk $m\in\Bbb N$, tetapkan

\Centerline{$\tilde f_m=\sum_{\sigma\in P',k_{\sigma}\le m}
\zeta_{\sigma}\innerprod{f}{\phi_{\sigma}}\phi_{\sigma}$.}

\noindent Maka, kapan pun $x$, $x'\in\Bbb R$ dan $|x-x'|\le 2^{-m}$, $|\tilde
f_m(x)|\le\bover12C_1\tilde f^*(x')$, di mana

\Centerline{$\tilde f^*(x') =\sup_{a\le x'\le
b,a<b}\Bover1{b-a}\int_a^b|\tilde f|$}

\noindent seperti dalam 286A. \Prf \ {\bf (i) } Karena $k_{\sigma}\ge
k_{\tau}$ untuk setiap $\sigma\in P'$, kita dapat menganggap bahwa $m\ge
k_{\tau}$. Misalkan $\hat J$ ialah interval diadik berpanjang $2^m$ yang memuat
$J_{\tau}$, dan $\hat y$ titik tengahnya. Tentukan $\psi=S_{-\hat
y}D_{2^{-m}/3}\varhat{\phi}$, yaitu
$\psi(y)=\varhat{\phi}(\bover132^{-m}(y-\hat y))$ untuk $y\in\Bbb R$.

\medskip

\quad{\bf (ii) } Jika $\sigma\in P'$ dan $k_{\sigma}\le m$ dan
$\varhat{\phi}_{\sigma}(y)\ne 0$, maka $y\in J^l_{\sigma}$. Tapi
$J_{\sigma}\cap\hat J\supseteq J_{\tau}$ tidak kosong, jadi
$J_{\sigma}\subseteq\hat J$, $|y-\hat y|\le\bover122^m$,
$|\bover132^{-m}(y-\hat y)|\le\bover16$ dan $\psi(y)=1$.

\medskip

\quad{\bf (iii) } Jika $\sigma\in P'$ dan $k_{\sigma}>m$ dan
$\varhat{\phi}_{\sigma}(y)\ne 0$, maka $J^r_{\sigma}\cap\hat J\supseteq
J^r_{\tau}$ tidak kosong, jadi $\hat J\subseteq J^r_{\sigma}$ dan $y\le
y_{\sigma}\le\hat y$; sekarang

\Centerline{$\hat y-y =(\hat y-y_{\sigma})+(y_{\sigma}-y) \ge\Bover12\cdot
2^m+\Bover1{20}\mu J_{\sigma} \ge\Bover35\cdot 2^m$, \quad$|\Bover13\cdot
2^{-m}(y-\hat y)|\ge\Bover15$}

\noindent dan $\psi(y)=0$. \medskip

\quad{\bf (iv)} Ini berarti bahwa jika $\sigma\in P'$ maka

$$\eqalign{\varhat{\phi}_{\sigma}\times\psi
&=\varhat{\phi}_{\sigma}\text{jika}k_{\sigma}\le m,\cr
&=0\text{jika}k_{\sigma}>m,\cr}$$

\noindent sehingga $\varhat{\vartildef}_m=\psi\times\varhat{\vartildef}$.

\medskip

\quad{\bf (v)} Dengan 283M, $\tilde f_m=\Bover1{\sqrt{2\pi}}\tilde
f*\varcheck{\psi}$, di mana $\tilde f*\varcheck{\psi}$ adalah konvolusi dari
$\tilde f$ dan transformasi Fourier invers $\varcheck{\psi}$ dari $\psi$.
(Secara ketat, 283M, dengan bantuan 284C, memberi tahu kita bahwa $\tilde
f_m$ dan $\Bover1{\sqrt{2\pi}}\tilde f*\varcheck{\psi}$ memiliki transformasi
Fourier yang sama. Menurut 283G, keduanya sama hampir di mana-mana; menurut
255K, konvolusi terdefinisi di mana-mana dan kontinu; jadi keduanya sebenarnya
merupakan fungsi yang sama di mana-mana.)

\Centerline{$\varcheck{\psi} =3\cdot 2^mM_{\hat y}D_{3\cdot 2^m}
\varhat{\phi}\varcheck{\phantom{\phi}} =3\cdot 2^mM_{\hat y}D_{3\cdot
2^m}\phi$,}

\noindent yaitu,

\Centerline{$\varcheck{\psi}(x) =3\cdot 2^me^{ix\hat y}\phi(3\cdot 2^mx)$}

\noindent untuk $x\in\Bbb R$.

\medskip

\quad{\bf (vi)} Tetapkan $w_1(x)=\min(w(3),w(x))$ untuk $x\in\Bbb R$, sehingga
$w_1$ tidak menurun pada $\ocint{-\infty,-3}$, tidak meningkat pada
$\coint{3,\infty}$, dan konstan pada $[-3,3]$, dan $|\phi(x)|\le C_1w_1(x)$
untuk setiap $x$, dengan memilih $C_1$ (286Ge). Ambil $x$, $x'\in\Bbb R$
sehingga $|x-x'|\le 2^{-m}$. Kemudian

$$\eqalignno{|\tilde f_m(x)|
&\le\Bover1{\sqrt{2\pi}}\int_{-\infty}^{\infty}
  |\tilde f(x-t)||\varcheck{\psi}(t)|dt
=\Bover{3\cdot 2^m}{\sqrt{2\pi}}\int_{-\infty}^{\infty}
  |\tilde f(x-t)||\phi(3\cdot 2^mt)|dt\cr
&\le\Bover{3\cdot 2^m}{\sqrt{2\pi}}C_1\int_{-\infty}^{\infty}
  |\tilde f(x-t)|w_1(3\cdot 2^mt)dt
=\Bover{3\cdot 2^m}{\sqrt{2\pi}}C_1\int_{-\infty}^{\infty}
  |\tilde f(x+t)|w_1(3\cdot 2^mt)dt\cr
\displaycause{karena $w_1$ merupakan fungsi genap}
&\le\Bover{3\cdot 2^m}{\sqrt{2\pi}}C_1
  \int_{-\infty}^{\infty}w_1(3\cdot 2^mt)dt
  \cdot\sup_{\Atop{a\le-2^{-m}}{b\ge 2^{-m}}}
  \Bover1{b-a}\int_a^b|\tilde f(x+t)|dt\cr
\displaycause{menurut 286B, karena $t\mapsto w_1(3\cdot 2^mt)$ tak menurun pada
$\ocint{-\infty,-2^{-m}}$, tak meningkat pada $\coint{2^{-m},\infty}$, dan
konstan pada $[-2^{-m},2^{-m}]$}
&=\Bover1{\sqrt{2\pi}}C_1\int_{-\infty}^{\infty}w_1
  \cdot\sup_{\Atop{a\le x-2^{-m}}{b\ge x+2^{-m}}}
  \Bover1{b-a}\int_a^b|\tilde f|
\le\Bover12C_1\int_{-\infty}^{\infty}w
  \cdot\tilde f^*(x')\cr
\displaycause{karena jika $a\le x-2^{-m}$ dan $b\ge x+2^{-m}$ maka
$a\le x'\le b$}
&=\Bover12C_1\tilde f^*(x')\cr}$$

\noindent (286Ga), sesuai kebutuhan.\ \Qed

\medskip

{\bf (i)} Untuk $x\in\Bbb R$, tetapkan

\Centerline{$v_3(x)=\bigl|\sum_{(\sigma,K)\in W_3}
\zeta_{\sigma}\innerprod{f}{\phi_{\sigma}}\phi_{\sigma}(x) \chi(E\cap
g^{-1}[J^r_{\sigma}]\cap K)(x)\bigr|$.}

\noindent Maka, kapan pun $L\in\Cal K$ dan $x$, $x'\in L$, $|v_3(x)|\le
C_1\tilde f^*(x')$. \Prf\ Kita bisa mengasumsikan bahwa $v_3(x)\ne 0$,
sehingga, khususnya, $x\in E$. Pasangan $(\sigma,K)$ yang menyumbang pada
jumlah yang membentuk $v_3(x)$ hanyalah pasangan dengan $x\in K$, sehingga
$K=L$, dan dengan $g(x)\in J^r_{\sigma}$. Selain itu, karena kita hanya meninjau
$\sigma\in T_{\tau}$, kita mempunyai $J^r_{\tau}\subseteq J^r_{\sigma}$.
Semua $J^r_{\sigma}$ merupakan interval diadik bersarang, berpanjang
$2^{k_{\sigma}-1}$, yang memuat $J^r_{\tau}$; karena itu ada $m$ sedemikian
sehingga (untuk $\sigma\in T_{\tau}$) $g(x)\in J^r_{\sigma}$ jika dan hanya
jika $k_{\sigma}\ge m$. Oleh karena itu

\Centerline{$v_3(x) =\bigl|\sum_{\sigma\in P',m\le k_{\sigma}<l_L}
\zeta_{\sigma}\innerprod{f}{\phi_{\sigma}}\phi_{\sigma}(x)\bigr| =|\tilde
f_{l_L-1}(x)-\tilde f_{m-1}(x)|$}

\noindent (kita harus mempunyai $m<l_L$ karena $v_3(x)\ne 0$). Sekarang
$|x-x'|\le 2^{-l_L}\le 2^{-m}$, sehingga (h) memberi tahu kita bahwa
$|\tilde f_{l_L-1}(x)|$ dan $|\tilde f_{m-1}(x)|$ masing-masing paling besar
$\bover12C_1\tilde f^*(x')$, dan $v_3(x)\le C_1\tilde f^*(x')$, seperti yang
diklaim.\ \Qed

Ini berarti bahwa $v_3(x)\le\Bover{C_1}{\mu L}\biggerint_L\tilde f^*$ untuk
setiap $x\in L$.

\wheader{286L}{6}{2}{2}{108pt}

{\bf (j)} Sekarang kita dalam posisi untuk memperkirakan

$$\eqalignno{|\alpha_3|
&=|\sum_{(\sigma,K)\in W_3}\zeta_{\sigma}\alpha_{\sigma K}|
\le\int_{-\infty}^{\infty}v_3
\le\sum_{\Atop{K\in\Cal K}{\mu K<\mu I_{\tau}}}\int_{G_K}v_3\cr
\displaycause{karena jika $v_3(x)\ne 0$ terdapat $(\sigma,K)\in W_3$
sedemikian sehingga $x\in K$, dan dalam hal ini $x\in G_K$ serta
$\mu K<\mu I_{\sigma}\le\mu I_{\tau}$}
&\le\sum_{\Atop{K\in\Cal K}{\mu K<\mu I_{\tau}}}
  \mu G_K\cdot\Bover{C_1}{\mu K}\int_K\tilde f^*\cr
\displaycause{menurut (i) di atas, karena $G_K\subseteq K$}
&\le C_1\sum_{\Atop{K\in\Cal K}{\mu K<\mu I_{\tau}}}
  \bover{2\gamma'}{w(\bover32)}\int_K\tilde f^*\cr
\displaycause{menurut (e)}
&\le\bover{2C_1\gamma'}{w(\bover32)}\int_{\hat I}\tilde f^*\cr
\displaycause{karena jika $\mu K<\mu I_{\tau}$ maka $K\subseteq\hat I$,
seperti dicatat dalam (a-iii)}
&\le\bover{2C_1\gamma'}{w(\bover32)}\sqrt{\mu\hat I}
  \cdot\|\tilde f^*\|_2\cr
\displaycause{menurut ketaksamaan Cauchy}
&\le\bover{2C_1\gamma'}{w(\bover32)}\sqrt{7\mu I_{\tau}}
  \cdot\sqrt{8}\|\tilde f\|_2\cr
\displaycause{menurut Teorema Maksimal, 286A}
&\le\bover{4C_1\gamma'\sqrt{14}}{w(\bover32)}\sqrt{\mu I_{\tau}}
  \cdot\gamma\sqrt{C_3\mu I_{\tau}}\cr
\displaycause{menurut (g)}
&=\bover{4C_1\sqrt{14 C_3}}{w(\bover32)}
  \gamma\gamma'\mu I_{\tau}.\cr}$$

\medskip

{\bf (k)} Mengumpulkan ini,

$$\eqalign{\sum_{\sigma\in P}\bigl|\innerprod{f}{\phi_{\sigma}}
  \int_{E\cap g^{-1}[J^r_{\sigma}]}\phi_{\sigma}\bigr|
&=\sum_{\Atop{\sigma\in P}{K\in\Cal K}}\zeta_{\sigma}\alpha_{\sigma K}
=\sum_{j=0}^3\sum_{(\sigma,K)\in W_j}\zeta_{\sigma}\alpha_{\sigma K}
\le\sum_{j=0}^3|\alpha_j|\cr
&\le\Bover72\cdot C_1\gamma\gamma'\mu I_{\tau}
  +\Bover87\cdot C_1\gamma\gamma'\mu I_{\tau}
  +28\cdot C_1\gamma\gamma'\mu I_{\tau}/w(\hbox{$\bover32$})\cr
&\qquad\qquad\qquad\qquad
  +4\sqrt{14C_3}\cdot
      C_1\gamma\gamma'\mu I_{\tau}/w(\hbox{$\bover32$})\cr
&=C_7\gamma\gamma'\mu I_{\tau},\cr}$$

\noindent seperti yang diklaim. }%end of proof of 286L

\leader{286M}{Lema Lacey--Thiele} Tetapkan $C_8=3C_7(C_5+C_6)$. Maka

\Centerline{$\sum_{\sigma\in Q}|\innerprod{f}{\phi_{\sigma}} \int_{E\cap
g^{-1}[J^r_{\sigma}]}\phi_{\sigma}| \le C_8$}

\noindent setiap kali $f\in\eusm L^2_{\Bbb C}$, $\|f\|_2=1$, $\mu E\le 1$, dan
$g:\Bbb R\to\Bbb R$ adalah terukur.

\proof{%
{\bf (a)} Langkah pertama adalah menggabungkan 286J dan 286K sebagai berikut:

jika $P\subseteq Q$ berhingga dan $\max(\sqrt{\mass_{Eg}(P)},\energy_f(P))\penalty-100\le\gamma$, terdapat $R\subseteq Q$ sehingga $\gamma^2\sum_{\tau\in R}\mu I_{\tau}\le C_5+C_6$ dan $\max(\sqrt{\mass_{Eg}(P\setminus R^+)},
\energy_f(P\setminus R^+))
\le\bover12\gamma$. \Prf\ Karena $\mass_{Eg}(P)\le\gamma^2$, 286J memberi
tahu kita bahwa terdapat $R_0\subseteq Q$ sehingga $\gamma^2\sum_{\tau\in R_0}\mu I_{\tau}\le C_5$ dan $\mass_{Eg}(P\setminus R_0^+)\le\bover14\gamma^2$. Beralih ke 286K, karena $\energy_f(P\setminus R_0^+)\le\energy_f(P)\le\gamma$, kita dapat menemukan $R_1\subseteq Q$ sehingga $\gamma^2\sum_{\tau\in R_1}\mu I_{\tau}\le C_6$ dan $\energy_f((P\setminus R_0^+)\setminus R_1^+)\le\bover12\gamma$. Tetapkan $R=R_0\cup R_1$.

\Centerline{$\gamma^2\sum_{\tau\in R}\mu I_{\tau}\le C_5+C_6$,
\quad$\mass_{Eg}(P\setminus R^+)\le\mass_{Eg}(P\setminus R_0^+)
\le\bover14\gamma^2$}

\noindent jadi $\max(\sqrt{\mass_{Eg}(P\setminus R^+)},\energy_f(P\setminus
R^+)) \le\bover12\gamma$.\ \Qed \fontdimen3\tenrm=1.67pt

\medskip

{\bf (b)} Sekarang ambil sembarang $P\subseteq Q$ berhingga. Pilih $k\in\Bbb N$
sedemikian sehingga $\max(\sqrt{\mass_{Eg}(P)},\energy_f(P))\le
2^k$. Dengan (a), kita dapat memilih $\sequencen{P_n}$, $\sequencen{R_n}$
induktif sehingga $P_0=P$ dan, untuk masing-masing $n\in\Bbb N$,

\Centerline{$P_{n+1}=P_n\setminus R_n^+$,}

\Centerline{$2^{2k-2n}\sum_{\tau\in R_n}\mu I_{\tau}\le C_5+C_6$,
\quad$\max(\sqrt{\mass_{Eg}(P_n)},\energy_f(P_n))\le 2^{k-n}$.}

\noindent Karena $\energy_f(\{\sigma\})=\sqrt{\mu
J_{\sigma}}|\innerprod{f}{\phi_{\sigma}}|>0$ setiap kali
$\innerprod{f}{\phi_{\sigma}}\ne 0$ (286H), $\innerprod{f}{\phi_{\sigma}}=0$
setiap kali $\sigma\in\bigcap_{n\in\Bbb N}P_n$, dan

$$\eqalignno{\sum_{\sigma\in P}\bigl|\innerprod{f}{\phi_{\sigma}}
  \int_{E\cap g^{-1}[J^r_{\sigma}]}\phi_{\sigma}\bigr|
&=\sum_{\sigma\in\bigcup_{n\in\Bbb N}P_n\setminus P_{n+1}}
  \bigl|\innerprod{f}{\phi_{\sigma}}
  \int_{E\cap g^{-1}[J^r_{\sigma}]}\phi_{\sigma}\bigr|\cr
&=\sum_{n=0}^{\infty}\sum_{\sigma\in P_n\setminus P_{n+1}}
  \bigl|\innerprod{f}{\phi_{\sigma}}
  \int_{E\cap g^{-1}[J^r_{\sigma}]}\phi_{\sigma}\bigr|\cr
&\le\sum_{n=0}^{\infty}\sum_{\tau\in R_n}
  \sum_{\Atop{\sigma\in P_n}{\sigma\ge\tau}}
     \bigl|\innerprod{f}{\phi_{\sigma}}
       \int_{E\cap g^{-1}[J^r_{\sigma}]}\phi_{\sigma}\bigr|\cr
&\le\sum_{n=0}^{\infty}\sum_{\tau\in R_n}C_7
  \energy_f(P_n)\mass_{Eg}(P_n)\mu I_{\tau}\cr
\displaycause{menurut 286L}
&\le C_7\sum_{n=0}^{\infty}2^{k-n}\min(1,2^{2k-2n})
  \sum_{\tau\in R_n}\mu I_{\tau}\cr
\displaycause{karena $\mass_{Eg}(P_n)\le 1$ untuk setiap $n$, seperti dicatat
dalam 286H}
&\le C_7\sum_{n=0}^{\infty}2^{k-n}\min(1,2^{2k-2n})2^{2n-2k}(C_5+C_6)
  \cr
&=C_7(C_5+C_6)\sum_{n=0}^{\infty}\min(2^{n-k},2^{k-n})\cr
&\le C_7(C_5+C_6)\sum_{n=-\infty}^{\infty}\min(2^n,2^{-n})
=3C_7(C_5+C_6).\cr}$$

\medskip

{\bf (c)} Karena ini benar untuk semua hingga $P\subseteq Q$,

\Centerline{$\sum_{\sigma\in Q}|\innerprod{f}{\phi_{\sigma}} \int_{E\cap
g^{-1}[J^r_{\sigma}]}\phi_{\sigma}| \le 3C_7(C_5+C_6)=C_8$,}

\noindent seperti yang diklaim. }%end of proof of 286M

\leader{286N}{Lema} Tetapkan $C_9=C_8\sqrt{2}$. Misalkan $f\in\eusm L^2_{\Bbb C}$,
$g:\Bbb R\to\Bbb R$ terukur, dan $\mu F<\infty$. Maka

\Centerline{$\sum_{\sigma\in Q}|\innerprod{f}{\phi_{\sigma}} \int_{F\cap
g^{-1}[J^r_{\sigma}]}\phi_{\sigma}| \le C_9\|f\|_2\sqrt{\mu F}$.}

\proof{ Hasilnya langsung berlaku jika $\|f\|_2=0$, yaitu jika $f=0$. Jadi kita
dapat menganggap $\|f\|_2>0$. Dengan membagi kedua ruas oleh $\|f\|_2$, kita
dapat menganggap $\|f\|_2=1$.

Pilih $k\in\Bbb Z$ sehingga $2^{k-1}<\mu F\le 2^k$. Kita mempunyai
permutasi $\sigma\mapsto\sigma^*:Q\to Q$ yang didefinisikan dengan menetapkan
bahwa $\sigma^*=(2^{-k}I_{\sigma},2^kJ_{\sigma})$; sehingga
$k_{\sigma^*}=k_{\sigma}+k$, $x_{\sigma^*}=2^{-k}x_{\sigma}$,
$y^l_{\sigma^*}=2^ky^l_{\sigma}$, $J^r_{\sigma^*}=2^kJ^r_{\sigma}$, dan untuk
setiap $x\in\Bbb R$

$$\eqalign{\phi_{\sigma}(2^kx)
&=\sqrt{\mu J_{\sigma}}e^{2^kiy^l_{\sigma}x}
  \phi((2^kx-x_{\sigma})\mu J_{\sigma})\cr
&=\sqrt{\mu J_{\sigma}}e^{iy^l_{\sigma^*}x}
  \phi(2^{k_{\sigma}+k}(x-2^{-k}x_{\sigma}))\cr
&=2^{-k/2}\sqrt{\mu J_{\sigma^*}}e^{iy^l_{\sigma^*}x}
  \phi((x-x_{\sigma^*})\mu J_{\sigma^*})
=2^{-k/2}\phi_{\sigma^*}(x).\cr}$$

\noindent Tuliskan $\tilde F=2^{-k}F$, sehingga $\mu\tilde F\le 1$, dan
$\tilde g(x)=2^kg(2^kx)$ untuk setiap $x$. Maka, untuk $\sigma\in Q$,

$$\eqalign{F\cap g^{-1}[J^r_{\sigma}]
&=\{x:x\in F,\,g(x)\in J^r_{\sigma}\}
=\{x:2^{-k}x\in\tilde F,\,2^{-k}\tilde g(2^{-k}x)\in J^r_{\sigma}\}\cr
&=\{x:2^{-k}x\in\tilde F,\,\tilde g(2^{-k}x)\in J^r_{\sigma^*}\}
=2^k\{x:x\in\tilde F,\,\tilde g(x)\in J^r_{\sigma^*}\}.\cr}$$

\noindent Tuliskan $\tilde f(x)=2^{k/2}f(2^kx)$, sehingga

\Centerline{$\|\tilde f\|_2=2^{k/2}\|D_{2^k}f\|_2=\|f\|_2=1$,}

\noindent sementara

\Centerline{$\innerprod{f}{\phi_{\sigma}}
=\int_{-\infty}^{\infty}f\times\bar\phi_{\sigma} =2^k\int_{-\infty}^{\infty}
f(2^kx)\overline{\phi_{\sigma}(2^kx)}dx =\innerprod{\tilde
f}{\phi_{\sigma^*}}$}

\noindent untuk setiap $\sigma\in Q$.

$$\eqalignno{\sum_{\sigma\in Q}\bigl|\innerprod{f}{\phi_{\sigma}}
  \int_{F\cap g^{-1}[J^r_{\sigma}]}\phi_{\sigma}\bigr|
&=2^k\sum_{\sigma\in Q}\bigl|\innerprod{f}{\phi_{\sigma}}
  \int_{2^{-k}(F\cap g^{-1}[J^r_{\sigma}])}
     \phi_{\sigma}(2^kx)dx\bigr|\cr
&=2^{k/2}\sum_{\sigma\in Q}\bigl|\innerprod{\tilde f}{\phi_{\sigma^*}}
  \int_{\tilde F\cap\tilde g^{-1}[J^r_{\sigma^*}]}
    \phi_{\sigma^*}\bigr|\cr
&=2^{k/2}\sum_{\tau\in Q}\bigl|\innerprod{\tilde f}{\phi_{\tau}}
  \int_{\tilde F\cap\tilde g^{-1}[J^r_{\tau}]}\phi_{\tau}\bigr|\cr
&\le 2^{k/2}C_8\cr
\displaycause{menurut lema Lacey--Thiele, yang diterapkan pada $\tilde g$,
$\tilde F$, dan $\tilde f$}
&\le C_9\sqrt{\mu F}
=C_9\|f\|_2\sqrt{\mu F}.\cr}$$
}%end of proof of 286N

\leader{286O}{Lema}\dvArevised{2013} (a) Untuk $z\in\Bbb R$, definisikan
$\theta_z:\Bbb R\to[0,1]$ dengan menetapkan

\Centerline{$\theta_z(y)=\varhat{\phi}(2^{-k}(y-\hat y))^2$}

\noindent setiap kali terdapat interval diadik $J\in\Cal I$ berpanjang $2^k$
sedemikian sehingga $z$ berada pada paruh kanan $J$, $y$ berada pada paruh
kiri $J$, dan $\hat y$ adalah titik seperempat bawah $J$; tetapkan nilainya nol
jika tidak ada $J$ demikian. Maka $(y,z)\mapsto\theta_z(y)$ terukur Borel,
$0\le\theta_z(y)\le
1$ untuk semua $y$, $z\in\Bbb R$, dan $\theta_z(y)=0$ jika $y\ge z$.

(b) Untuk $k\in\Bbb Z$, tetapkan $Q_k=\{\sigma:\sigma\in Q$, $k_{\sigma}=k\}$.
Misalkan $[Q]^{<\omega}$ ialah himpunan semua bagian berhingga dari $Q$,
$[\Bbb Z]^{<\omega}$ himpunan semua bagian berhingga dari $\Bbb Z$, dan $\Cal
L$ himpunan semua bagian $L$ dari $Q$ sedemikian sehingga $L\cap Q_k$
berhingga untuk setiap $k$. Jika $K\in[\Bbb Z]^{<\omega}$ dan $L\in\Cal
L$, himpunan

$$\eqalign{\Cal P_{KL}
&=\{P:P\in[Q]^{<\omega},\,P\cap Q_k\supseteq L\cap Q_k
  \text{setiap kali }k\in\Bbb Z
\cr&\mskip200mu  
   \text{ dan }k\in K\text{ atau }P\cap Q_k\ne\emptyset\};\cr}$$
   
\noindent tetapkan

\Centerline{$\Cal F=\{\Cal P:\Cal P\subseteq[Q]^{<\omega}$ dan ada $K\in[\Bbb
Z]^{<\omega}$, $L\in\Cal L$ sehingga $\Cal P\supseteq\Cal P_{KL}\}$.}
  
\noindent Maka $\Cal F$ merupakan filter pada $[Q]^{<\omega}$ dan

\Centerline{$2\pi\int_F(\varhat h\times\theta_z)\varspcheck =\lim_{P\to\Cal
F}\sum_{\sigma\in P,z\in J^r_{\sigma}}
\innerprod{h}{\phi_{\sigma}}\int_F\phi_{\sigma}$}
   
\noindent untuk setiap $z\in\Bbb R$, setiap fungsi uji yang menurun cepat $h$,
dan setiap himpunan terukur $F\subseteq\Bbb R$ berukuran hingga.

\proof{{\bf (a)(i) } Mula-mula kita jelaskan mengapa aturan di atas
mendefinisikan fungsi $\theta_z$. Misalkan $M$ ialah himpunan semua $k\in\Bbb Z$
sedemikian sehingga $z$ berada pada paruh kanan interval diadik $\hat J_k$
berpanjang $2^k$ yang memuat $z$. Untuk $k\in M$, misalkan $\hat y_k$ ialah
titik tengah paruh kiri $\hat J^l_k$ dari $\hat J_k$, dan tetapkan
$\psi_k(y)=\varhat{\phi}(2^{-k}(y-\hat y_k))^2$ untuk $y\in\Bbb R$; maka
$\psi_k$ mulus dan nol di luar $\hat J^l_k$. Perhatikan bahwa
bahwa jika $k$, $k'$ adalah anggota yang berbeda dari $M$, maka $\hat J^l_k$
dan $\hat J^l_{k'}$ saling lepas, seperti dicatat dalam 286E(b-iv). Jadi
$\theta_z$ hanyalah jumlah $\sum_{k\in M}\psi_k$. Karena $\varhat{\phi}$
mengambil nilai dalam $[0,1]$, maka juga $\theta_z$. Jika $y\ge z$, maka tentu
saja $y\notin\hat J^l_k$ untuk setiap $k\in M$, jadi $\theta_z(y)=0$.

\medskip

\quad{\bf (ii)} Untuk melihat bahwa $(y,z)\mapsto\theta_z(y)$ terukur Borel,
perhatikan bahwa

\Centerline{$\{(y,z):\theta_z(y)\ge\gamma\} =\bigcup_{\sigma\in Q}\{(y,z):
\hat\phi((y-y^l_{\sigma})\mu I_{\sigma})^2\ge\gamma$, $z\in J^r_{\sigma}\}$}

\noindent untuk setiap $\gamma\in\Bbb R$.

\medskip

{\bf (b)(i) } $\emptyset$ termasuk dalam $[\Bbb Z]^{<\omega}$ dan $\Cal L$,
dan $[Q]^{<\omega}=\Cal P_{\emptyset\emptyset}$ termasuk dalam $\Cal F$. Jika
$K\in[\Bbb Z]^{<\omega}$ dan $L\in\Cal L$, maka $\bigcup_{k\in K}L\cap Q_k$
termasuk dalam $\Cal P_{KL}$. Jadi tidak ada $\Cal P_{KL}$ yang kosong dan
$\emptyset\notin\Cal F$.

Jika $\Cal P$, $\Cal P'\in\Cal F$, ada $K$, $K'\in[\Bbb Z]^{<\omega}$ dan $L$,
$L'\in\Cal L$ sehingga $\Cal P_{KL}\subseteq\Cal P$ dan $\Cal
P_{K'L'}\subseteq\Cal P'$. Sekarang $K\cup K'\in[\Bbb Z]^{<\omega}$, $L\cup
L'\in\Cal L$ dan

\Centerline{$\Cal P_{K\cup K',L\cup L'} \subseteq\Cal P_{KL}\cap\Cal
P_{K'L'}\subseteq\Cal P\cap\Cal P'$,}

\noindent jadi $\Cal P\cap\Cal P'\in\Cal F$.

Jika $\Cal P\in\Cal F$ dan $\Cal P\subseteq\Cal P'\subseteq[Q]^{<\omega}$,
maka tentu saja $\Cal P'\in\Cal F$. Jadi $\Cal F$ adalah filter pada
$[Q]^{<\omega}$.

\medskip

\quad{\bf (ii)} Sekarang pilih $z\in\Bbb R$, fungsi uji yang menurun cepat
$h$, dan himpunan $F$ berukuran hingga. Ambil $M$ dan $\psi_k$,
$\hat J_k$, $\hat y_k$ untuk $k\in M$ dari (a-i) di atas; akan memudahkan jika
kita menetapkan $\psi_k=0$ untuk $k\in\Bbb Z\setminus M$, sehingga
$\theta_z=\sum_{k\in\Bbb Z}\psi_k$.

Untuk $k\in\Bbb Z$,

\Centerline{$2\pi\int_F(\varhat h\times\psi_k)\varspcheck =\sum_{\sigma\in
Q_k,z\in J^r_{\sigma}} \innerprod{h}{\phi_{\sigma}}\int_F\phi_{\sigma}$.}
  
\noindent\Prf\ Jika $k\notin M$, maka $z\notin J^r_{\sigma}$ untuk setiap
$\sigma\in Q_k$, sedangkan $\psi_k=0$, sehingga hasilnya langsung berlaku.
Jadi anggap $k\in M$ dan $\hat y_k$ terdefinisi.
$\sigma\in Q_k$ dan $z\in J^r_{\sigma}$, $y^l_{\sigma}=\hat y_k$ dan
$x_{\sigma}$ berbentuk $2^{-k}(n+\bover12)$ untuk suatu $n\in\Bbb Z$. Jadi

$$\eqalignno{\innerprod{h}{\phi_{\sigma}}
&=\int_{-\infty}^{\infty}
  \varhat h\times\bar\varhat{\phi}_{\sigma}\cr
\displaycause{284O}
&=\int_{-\infty}^{\infty}\varhat h(t)
  \cdot 2^{-k/2}e^{2^{-k}i(n+\bover12)(t-\hat y_k)}
  \varhat\phi(2^{-k}(t-\hat y_k))dt\cr
\displaycause{menurut rumus dalam 286Eb, karena $\varhat{\phi}$ bernilai riil}
&=2^{k/2}\int_{-\infty}^{\infty}\varhat h(2^kt+\hat y_k)
  e^{i(n+\bover12)t}\varhat\phi(t)dt
=2^{k/2}\int_{-\pi}^{\pi}\varhat h(2^kt+\hat y_k)
  e^{i(n+\bover12)t}\varhat\phi(t)dt\cr
\displaycause{karena $\varhat{\phi}(t)=0$ jika $|t|\ge\bover15$}
&=2^{k/2}\int_{-\pi}^{\pi}g(t)e^{int}dt,\cr}$$

\noindent dengan $g(t)=\varhat h(2^kt+\hat y_k)e^{it/2}\varhat{\phi}(t)$
untuk $-\pi<t\le\pi$. Jadi, jika kita menetapkan
$c_n=\Bover1{2\pi}\int_{-\pi}^{\pi}g(t)e^{-int}dt$, seperti di 282A, kita
punya

\Centerline{$\innerprod{h}{\phi_{\sigma}}=2^{k/2}\cdot 2\pi c_{-n}$}

\noindent ketika $\sigma\in Q_k$, $z\in J^r_{\sigma}$, dan
$x_{\sigma}=2^{-k}(n+\bover12)$. Karena $g$ mulus dan nol di
luar $[-\bover15,\bover15]$, $\sum_{n=-\infty}^{\infty}|c_n|<\infty$ (282Rb).

Sekarang, untuk setiap $y\in\hat J^l_k$, dengan menuliskan $R_k$ untuk

\Centerline{$\{\sigma:\sigma\in Q_k$, $z\in J^r_{\sigma}\} =\{\sigma:\sigma\in
Q_k$, $J_{\sigma}=\hat J_k\} =\{(I,\hat J_k):I\in\Cal I$, $\mu I=2^{-k}\}$,}

\noindent kita mempunyai

$$\eqalignno{\sum_{\sigma\in R_k}
  \innerprod{h}{\phi_{\sigma}}\varhat{\phi}_{\sigma}(y)
&=\sum_{n=-\infty}^{\infty}2^{k/2}\cdot 2\pi c_{-n}
   \cdot 2^{-k/2}e^{-2^{-k}i(n+\bover12)(y-\hat y_k)}
   \varhat{\phi}(2^{-k}(y-\hat y_k))\cr
&=2\pi\varhat{\phi}(2^{-k}(y-\hat y_k))
   e^{-2^{-k-1}i(y-\hat y_k)}
   \sum_{n=-\infty}^{\infty}c_{-n}e^{-2^{-k}in(y-\hat y_k)}\cr
&=2\pi\varhat{\phi}(2^{-k}(y-\hat y_k))
   e^{-2^{-k-1}i(y-\hat y_k)}
   \sum_{n=-\infty}^{\infty}c_ne^{in 2^{-k}(y-\hat y_k)}\cr
&=2\pi\varhat{\phi}(2^{-k}(y-\hat y_k))
   e^{-2^{-k-1}i(y-\hat y_k)}g(2^{-k}(y-\hat y_k))\cr
\displaycause{menurut 282L(i), karena $2^{-k}|y-\hat y_k|\le\bover14<\pi$
dan $g$ mulus}
&=2\pi e^{-2^{-k-1}i(y-\hat y_k)}\varhat{\phi}(2^{-k}(y-\hat y_k))
  \varhat h(y)e^{2^{-k-1}i(y-\hat y_k)}
  \varhat{\phi}(2^{-k}(y-\hat y_k))\cr
&=2\pi\varhat h(y)\psi_k(y).\cr}$$

\noindent Di sisi lain, jika $y\in\Bbb R\setminus\hat J^l_k$,
$\psi_k(y)=\varhat{\phi}_{\sigma}(y)=0$ untuk setiap $\sigma\in R_k$, jadi
sekali lagi $\sum_{\sigma\in R_k}
\innerprod{h}{\phi_{\sigma}}\varhat{\phi}_{\sigma}(y) =2\pi\varhat
h(y)\psi_k(y)$.

Selanjutnya,

\Centerline{$\sum_{\sigma\in R_k}|\innerprod{h}{\phi_{\sigma}}| =2\pi\cdot
2^{k/2}\sum_{n=-\infty}^{\infty}|c_n|$}

\noindent dan

\Centerline{$\sup_{\sigma\in R_k}\int_{-\infty}^{\infty}
|\varhat{\phi}_{\sigma}| =2^{k/2}\int_{-\infty}^{\infty}|\varhat{\phi}|$}

\noindent berhingga, sedangkan $\varhat{\chi F}$ tentu terbatas. Jadi

$$\eqalignno{2\pi\int_F(\varhat h\times\psi_k)\varspcheck
&=2\pi\innerprod{\varhat h\times\psi_k)\varspcheck}{\chi F}
=2\pi\innerprod{(\varhat h\times\psi_k)\varspcheck)\varsphat}
  {\varhat{\chi F}}\cr
\displaycause{sekali lagi menurut 284Ob}
&=2\pi\innerprod{\varhat h\times\psi_k}{\varhat{\chi F}}
=\int_{-\infty}^{\infty}
    2\pi\varhat h\times\psi_k\times\overline{\varhat{\chi F}}
\cr&=\int_{-\infty}^{\infty}\sum_{\sigma\in R_k}
   \innerprod{h}{\phi_{\sigma}}\varhat{\phi}_{\sigma}
   \times\overline{\varhat{\chi F}}
=\sum_{\sigma\in R_k}\innerprod{h}{\phi_{\sigma}}\int_{-\infty}^{\infty}
   \varhat{\phi}_{\sigma}\times\overline{\varhat{\chi F}}\cr
\displaycause{226E}   
&=\sum_{\sigma\in R_k}\innerprod{h}{\phi_{\sigma}}
   \int_F{\phi}_{\sigma}
=\sum_{\sigma\in Q_k,z\in J^r_{\sigma}}\innerprod{h}{\phi_{\sigma}}
   \int_F{\phi}_{\sigma}.  \text{\Qed}\cr}$$
   
\medskip

\quad{\bf (iii) } Dalam kalimat terakhir argumen di atas, kita menggunakan
teorema B. Levi dalam bentuk 226E, meskipun $R_k$ memiliki pengindeksan alami,
karena nanti kita secara khusus memerlukan pernyataan bahwa

\inset{\noindent untuk setiap $\epsilon>0$ ada hingga $L_0\subseteq R_k$
sedemikian sehingga

\centerline{$|2\pi\int_F(\varhat h\times\psi_k)\varspcheck -\sum_{\sigma\in
L}\innerprod{h}{\phi_{\sigma}} \int_F{\phi}_{\sigma}|\le\epsilon$}
   
\noindent setiap kali $L\subseteq R_k$ adalah hingga dan $L\supseteq L_0$;}

\noindent dari sini segera diperoleh bahwa
   
\inset{\noindent untuk setiap $\epsilon>0$ ada hingga $L_0\subseteq Q_k$
sedemikian sehingga

\centerline{$|2\pi\int_F(\varhat h\times\psi_k)\varspcheck -\sum_{\sigma\in
L,z\in J^r_{\sigma}}\innerprod{h}{\phi_{\sigma}}
\int_F{\phi}_{\sigma}|\le\epsilon$}
   
\noindent setiap kali $L\subseteq Q_k$ adalah hingga dan $L\supseteq L_0$.}

\medskip

\quad{\bf (iv)} Sekarang tinjau $(\varhat
h\times\theta_z)\varspcheck$. Karena setiap $\psi_k$ nonnegatif,
$\theta_z=\sum_{k\in\Bbb Z}\psi_k$ adalah terbatas di atas oleh $1$, dan
$\varhat h$ adalah terintegralkan,

$$\eqalign{\int_F(\varhat h\times\theta_z)\varspcheck
&=\int_{-\infty}^{\infty}\varhat h\times\theta_z\times\varhat{\chi F}
\cr&=\sum_{k\in\Bbb Z}\int_{-\infty}^{\infty}\varhat h\times\psi_k\times\varhat{\chi F}
=\sum_{k\in\Bbb Z}\int_F(\varhat h\times\psi_k)\varspcheck.\cr}$$

\noindent Jadi kita dapat mengatakan

\inset{\noindent untuk setiap $\epsilon>0$ ada $K_0\in[\Bbb Z]^{<\omega}$
sedemikian sehingga

\centerline{$|\int_F(\varhat h\times\theta_z)\varspcheck -\sum_{k\in
K}\int_F(\varhat h\times\psi_k)\varspcheck|\le\epsilon$}

\noindent setiap kali $K\in[\Bbb Z]^{<\omega}$ dan $K\supseteq K_0$.}

\medskip

\quad{\bf (v)} Untuk mengekspresikan fakta di atas dalam hal batas sepanjang
filter $\Cal F$, kita dapat berargumen sebagai berikut. Ambil $\epsilon>0$.
Untuk setiap $k\in\Bbb Z$, (iii) memberi tahu kita bahwa ada himpunan berhingga
$L_k\subseteq Q_k$ sedemikian sehingga

\Centerline{$|2\pi\int_F(\varhat h\times\psi_k)\varspcheck -\sum_{\sigma\in
L',z\in J^r_{\sigma}}\innerprod{h}{\phi_{\sigma}} \int_F{\phi}_{\sigma}|\le
2^{-|k|}\epsilon$}
   
\noindent setiap kali $L'\subseteq Q_k$ adalah hingga dan $L'\supseteq L_k$;
tentu saja kita dapat menganggap bahwa setiap $L_k$ tidak kosong. Tetapkan
$L=\bigcup_{k\in\Bbb Z}L_k$, sehingga $L\cap Q_k=L_k$ adalah hingga untuk
setiap $k$, dan $L\in\Cal L$. Selanjutnya, ada $K\in[\Bbb Z]^{<\omega}$
sedemikian sehingga

\Centerline{$|\int_F(\varhat h\times\theta_z)\varspcheck -\sum_{k\in
K'}\int_F(\varhat h\times\psi_k)\varspcheck| \le\epsilon$}

\noindent setiap kali $K'\in[\Bbb Z]^{<\omega}$ dan $K'\supseteq K$. Ambil
$P\in\Cal P_{KL}$. Menentukan $K'=\{k:P\cap Q_k\ne\emptyset\}$, kita memiliki
$K'\supseteq K$ dan $P\cap Q_k\supseteq L\cap Q_k$ untuk setiap $k\in K'$.

$$\eqalign{&|2\pi\int_F(\varhat h\times\theta_z)\varspcheck
   -\sum_{\sigma\in P,z\in J^r_{\sigma}}
      \innerprod{h}{\phi_{\sigma}}\int_F\phi_{\sigma}|
\cr&\mskip100mu
=|2\pi\int_F(\varhat h\times\theta_z)\varspcheck
   -\sum_{k\in K'}\sum_{\Atop{\sigma\in P\cap Q_k}{z\in J^r_{\sigma}}}
      \innerprod{h}{\phi_{\sigma}}\int_F\phi_{\sigma}|
\cr&\mskip100mu
\le 2\pi|\int_F(\varhat h\times\theta_z)\varspcheck
   -\sum_{k\in K'}\int_F(\varhat h\times\psi_k)\varspcheck|      
\cr&\mskip200mu
+\sum_{k\in K'}|2\pi\int_F(\varhat h\times\psi_k)\varspcheck
   -\sum_{\Atop{\sigma\in P\cap Q_k}{z\in J^r_{\sigma}}}
      \innerprod{h}{\phi_{\sigma}}\int_F\phi_{\sigma}|
\cr&\mskip100mu
\le 2\pi\epsilon+\sum_{k\in K'}2^{-|k|}\epsilon
\le (2\pi+3)\epsilon.\cr}$$

\noindent Karena $\Cal P_{KL}\in\Cal F$ dan $\epsilon$ sembarang,

\Centerline{$2\pi\int_F(\varhat h\times\theta_z)\varspcheck =\lim_{P\to\Cal
F}\sum_{\sigma\in P,z\in J^r_{\sigma}}
\innerprod{h}{\phi_{\sigma}}\int_F\phi_{\sigma}$}

\noindent seperti yang diklaim. }%end of proof of 286O
   
\leader{286P}{Lema}\dvArevised{2013} Misalkan $h$ fungsi uji yang menurun
cepat. Untuk $x\in\Bbb R$, tetapkan

\Centerline{$Ah(x) =\sup_{z\in\Bbb R}|2\pi(\varhat
h\times\theta_z)\varspcheck(x)|$.}

\noindent Maka $Ah:\Bbb R\to[0,\infty]$ adalah Borel terukur, dan $\int_FAh\le
4C_9\|h\|_2\sqrt{\mu F}$ kapan pun $\mu F<\infty$.

\proof{{\bf (a)} Karena $(\varhat h\times\theta_z)\varspcheck$ kontinu
untuk setiap $z$, $Ah$ semikontinu bawah, dan karena itu
Borel terukur, oleh 256Ma. Oleh 256Mb,

\Centerline{$\int_FAh =\sup\{\int_F\sup_{i\le n} |2\pi(\varhat
h\times\theta_{z_i})\varspcheck|: z_0,\ldots,z_n\in\Bbb R\}$.}

\medskip

{\bf (b)} Untuk sementara, pilih $z_0,\ldots,z_n\in\Bbb R$.

\medskip

\quad{\bf (i)} Tetapkan $v_i=2\pi(\varhat h\times\theta_{z_i})\varspcheck$ untuk
$i\le n$, dan $v=\sup_{i\le n}|v_i|$. Tetapkan
$E_i=\{x:v(x)=|v_i(x)|\}\setminus\bigcup_{j<i}\{x:v(x)=|v_j(x)|\}$ untuk $i\le
n$, sehingga $(E_0,\ldots,E_n)$ adalah partisi dari $\Bbb R$ menjadi
himpunan-himpunan Borel, dan

\Centerline{$\int_Fv =\int_F\sum_{i=0}^n|v_i|\times\chi E_i
=\int_F|\sum_{i=0}^nv_i\times\chi E_i| \le
4|\int_{F'}\sum_{i=0}^nv_i\times\chi E_i|$}

\noindent untuk suatu $F'\subseteq F$ terukur yang sesuai (246K). Dengan menetapkan
$g(x)=z_i$ untuk $x\in E_i$, $g:\Bbb R\to\Bbb R$ adalah Borel terukur.

\medskip

\quad{\bf (ii)} Untuk setiap $i\le n$,

$$\eqalignno{\int_{F'}v_i\times\chi E_i
&=\int_{F'\cap E_i}v_i
=\lim_{P\to\Cal F}\sum_{\sigma\in P,z_i\in J^r_{\sigma}}
  \innerprod{h}{\phi_{\sigma}}\int_{F'\cap E_i}\phi_{\sigma}\cr
\displaycause{dengan $\Cal F$ filter pada $[Q]^{<\omega}$ yang diuraikan dalam
286O}
&=\lim_{P\to\Cal F}\int_{F'\cap E_i}
  \sum_{\sigma\in P,z_i\in J^r_{\sigma}}
  \innerprod{h}{\phi_{\sigma}}\phi_{\sigma}
\cr&=\lim_{P\to\Cal F}\int_{F'\cap E_i}
  \sum_{\sigma\in P,g(x)\in J^r_{\sigma}}
  \innerprod{h}{\phi_{\sigma}}\phi_{\sigma}(x)dx.\cr}$$
  
\noindent Jadi

$$\eqalign{\int_{F'}\sum_{i=0}^nv_i\times\chi E_i  
&=\lim_{P\to\Cal F}\sum_{i=0}^n\int_{F'\cap E_i}
  \sum_{\sigma\in P,g(x)\in J^r_{\sigma}}
  \innerprod{h}{\phi_{\sigma}}\phi_{\sigma}(x)dx
\cr&=\lim_{P\to\Cal F}\int_{F'}
  \sum_{\sigma\in P,g(x)\in J^r_{\sigma}}
  \innerprod{h}{\phi_{\sigma}}\phi_{\sigma}(x)dx.\cr}$$
  
\noindent Sekarang, untuk setiap himpunan berhingga $P\subseteq Q$,

\Centerline{$\int_{F'}\sum_{\sigma\in P,g(x)\in J^r_{\sigma}}
\innerprod{h}{\phi_{\sigma}}\phi_{\sigma}(x)dx =\sum_{\sigma\in P}\int_{F'\cap
g^{-1}[J^r_{\sigma}]} \innerprod{h}{\phi_{\sigma}}\phi_{\sigma}$;}
  
\noindent Hal ini dapat dipandang sebagai penerapan teorema Fubini jika $Q$
diberi ukuran hitung dan kita meninjau fungsi

$$\eqalign{(x,\sigma)
&\mapsto\innerprod{h}{\phi_{\sigma}}\phi_{\sigma}(x)
  \text{ jika }x\in F',\,\sigma\in P\text{ dan }g(x)\in J^r_{\sigma},
\cr&\mapsto 0\text{ selainnya}.\cr}$$

\noindent Ini berarti bahwa

\Centerline{$|\int_{F'}\sum_{\sigma\in P,g(x)\in J^r_{\sigma}}
\innerprod{h}{\phi_{\sigma}}\phi_{\sigma}(x)dx| \le\sum_{\sigma\in
P}|\innerprod{h}{\phi_{\sigma}} \int_{F'\cap
g^{-1}[J^r_{\sigma}]}\phi_{\sigma}| \le C_9\|h\|_2\sqrt{\mu F'}$}

\noindent menurut 286N. Dengan mengambil limit saat $P\to\Cal F$,

\Centerline{$|\int_{F'}\sum_{i=0}^nv_i\times\chi E_i| \le C_9\|h\|_2\sqrt{\mu
F'}$.}

\medskip

\quad{\bf (iii)} Jadi kita memiliki

$$\eqalign{\int_F\sup_{i\le n}
   |2\pi(\varhat h\times\theta_{z_i})\varspcheck|
=\int_Fv
&\le 4|\int_{F'}\sum_{i=0}^nv_i\times\chi E_i|
\cr&\le 4C_9\|h\|_2\sqrt{\mu F'}
\le 4C_9\|h\|_2\sqrt{\mu F}.\cr}$$

\medskip

{\bf (c)} Karena $z_0,\ldots,z_n$ sembarang, $\int_FAh\le 4C_9\|h\|_2\sqrt{\mu F}$, seperti yang diklaim. }%end of proof of 286P

\leader{286Q}{Lema} Untuk $\alpha>0$ dan $y$, $z$, $\beta\in\Bbb R$,
tetapkan $\theta'_{z\alpha\beta}(y)=\theta_{\alpha z+\beta}(\alpha y+\beta)$.
Maka

(a) pemetaan $(\alpha,\beta,y,z)\mapsto\theta'_{z\alpha\beta}(y):
\ooint{0,\infty}\times\BbbR^3\to[0,1]$ adalah Borel terukur;

(b) $\theta'_{z\alpha\beta}(y)=0$ setiap kali $y\ge z$;

(c) untuk setiap fungsi uji yang menurun cepat $h$ dan setiap $z\in\Bbb R$,

\Centerline{$2\pi|(\varhat h\times\theta'_{z\alpha\beta})\varspcheck| \le
D_{1/\alpha}AM_{\beta}D_{\alpha}h$}

\noindent\cmmnt{(dalam notasi 286C)} di setiap titik.

\proof{{\bf (a)} Kita hanya perlu mengingat bahwa
$(y,z)\mapsto\theta_z(y):\BbbR^2\to\Bbb R$ adalah Borel terukur (286Oa), dan
bahwa $(\alpha,\beta,y,z)\mapsto\theta'_{z\alpha\beta}(y)$ dibangun dari ini,
$+$ dan $\times$.

\medskip

{\bf (b)} Sekali lagi, ini langsung dari 286Oa, karena $\alpha>0$.

\medskip

{\bf (c)} Tetapkan $v=\alpha z+\beta$, sehingga
$\theta'_{z\alpha\beta}=D_{\alpha}S_{\beta}\theta_v$. Maka

$$\eqalign{\varhat h\times\theta'_{z\alpha\beta}
&=\varhat h\times D_{\alpha}S_{\beta}\theta_v
=D_{\alpha}S_{\beta}(S_{-\beta}D_{1/\alpha}\varhat h\times\theta_v)\cr
&=\alpha D_{\alpha}S_{\beta}
  (S_{-\beta}(D_{\alpha}h)\varsphat\times\theta_v)
=\alpha D_{\alpha}S_{\beta}
  ((M_{\beta}D_{\alpha}h)\varsphat\times\theta_v),\cr}$$

\noindent jadi

$$\eqalign{(\varhat h\times\theta'_{z\alpha\beta})\varspcheck
&=\alpha\bigl(D_{\alpha}S_{\beta}
  ((M_{\beta}D_{\alpha}h)\varsphat\times\theta_v)\bigr)\varspcheck\cr
&=D_{1/\alpha}\bigl(S_{\beta}
  ((M_{\beta}D_{\alpha}h)\varsphat\times\theta_v)\bigr)\varspcheck
=D_{1/\alpha}M_{-\beta}\bigl(
  (M_{\beta}D_{\alpha}h)\varsphat\times\theta_v\bigr)\varspcheck\cr}$$

\noindent dan

\Centerline{$2\pi|(\varhat h\times\theta'_{z\alpha\beta})\varspcheck|
=2\pi D_{1/\alpha}\bigl|\bigl(
  (M_{\beta}D_{\alpha}h)\varsphat
  \times\theta_v\bigr)\varspcheck\bigr|
\le D_{1/\alpha}A(M_{\beta}D_{\alpha}h)$.}
}%end of proof of 286Q

\leader{286R}{Lema} Untuk setiap $y$, $z\in\Bbb R$,

\Centerline{$\tilde\theta_z(y) =\int_1^2\Bover1{\alpha}\bigl(\lim_{n\to\infty}
\Bover1n\int_0^n\theta'_{z\alpha\beta}(y)d\beta\bigr)d\alpha$}

\noindent didefinisikan, dan

$$\eqalign{\tilde\theta_z(y)&=\tilde\theta_1(0)>0\text{jika}y<z,\cr
&=0\text{jika}y\ge z.\cr}$$

\proof{{\bf (a)} Kasus $y\ge z$ langsung, karena jika $y\ge z$ maka
$\theta'_{z\alpha\beta}(y)=0$ untuk semua $\alpha>0$ dan $\beta\in\Bbb R$
(286Qb) dan $\tilde\theta_z(y)=0$. Jadi, untuk sisa bukti, kita meninjau kasus
$y<z$.

\medskip

{\bf (b)(i) } Jika $y<z\in\Bbb R$ dan $\alpha>0$, tetapkan
$l=\lfloor\log_2(20\alpha(z-y))\rfloor$. Kemudian
$\theta'_{z,\alpha,\beta+2^l}(y)=\theta'_{z\alpha\beta}(y)$ untuk setiap
$\beta\in\Bbb R$. \Prf\ Jika $\theta'_{z\alpha\beta}(y)=\theta_{\alpha
z+\beta}(\alpha y+\beta)$ tidak nol, harus ada $k$, $m\in\Bbb Z$ sedemikian
sehingga

\Centerline{$2^k(m+\Bover12)\le\alpha z+\beta<2^k(m+1)$}

\noindent dan

\Centerline{$\varhat{\phi}(2^{-k}(\alpha y+\beta)-(m+\Bover14))^2
=\theta'_{z\alpha\beta}(y)\ne 0$,}

\noindent jadi

\Centerline{$2^km\le\alpha y+\beta\le 2^k(m+\Bover9{20})$}

\noindent karena $\varhat\phi$ adalah nol di luar $[-\bover15,\bover15]$.
Dalam hal ini, $\Bover1{20}\cdot 2^k<\alpha(z-y)$, sehingga $k\le l$.

\Centerline{$2^k(m+2^{l-k}+\Bover12)\le\alpha z+\beta+2^l <2^k(m+2^{l-k}+1)$,}

\Centerline{$2^k(m+2^{l-k})\le\alpha y+\beta+2^l <2^k(m+2^{l-k}+\Bover12)$,}

\noindent jadi

\Centerline{$\theta'_{z,\alpha,\beta+2^l}(y) =\varhat{\phi}(2^{-k}(\alpha
y+\beta+2^l)-(m+2^{l-k}+\Bover14))^2 =\theta'_{z\alpha\beta}(y)$.}

\noindent Sama halnya,

\Centerline{$2^k(m-2^{l-k}+\Bover12)\le\alpha z+\beta-2^l <2^k(m-2^{l-k}+1)$,}

\Centerline{$2^k(m-2^{l-k})\le\alpha y+\beta-2^l <2^k(m-2^{l-k}+\Bover12)$,}

\noindent jadi

\Centerline{$\theta'_{z,\alpha,\beta-2^l}(y) =\varhat{\phi}(2^{-k}(\alpha
y+\beta-2^l)-(m-2^{l-k}+\Bover14))^2 =\theta'_{z\alpha\beta}(y)$.}

\noindent Ini menunjukkan bahwa
$\theta'_{z,\alpha,\beta+2^l}(y)=\theta'_{z\alpha\beta}(y)$ jika salah satunya
tidak nol; karena itu kesamaan berlaku dalam semua kasus.\ \Qed

\medskip

{\bf (ii)} Akibatnya $g(\alpha,y,z)
=\lim_{b\to\infty}\Bover1b\biggerint_0^b\theta'_{z\alpha\beta}(y)d\beta$
terdefinisi. \Prf\ Tetapkan

\Centerline{$\gamma =2^{-l}\int_0^{2^l}\theta'_{z\alpha\beta}(y)d\beta$.}

\noindent Dari (i) kita melihat bahwa

\Centerline{$\gamma
=2^{-l}\int_{2^lm}^{2^l(m+1)}\theta'_{z\alpha\beta}(y)d\beta$}

\noindent untuk setiap $m\in\Bbb Z$, dan oleh karena itu bahwa

\Centerline{$\gamma
=\Bover1{2^lm}\int_0^{2^lm}\theta'_{z\alpha\beta}(y)d\beta$}

\noindent untuk setiap $m\ge 1$. Karena $\theta'_{z\alpha\beta}(y)$ selalu
lebih besar atau sama dengan $0$, jika $2^lm\le b\le 2^l(m+1)$ maka

\Centerline{$\Bover{m}{m+1}\gamma
=\Bover1{2^l(m+1)}\int_0^{2^lm}\theta'_{z\alpha\beta}
\le\Bover1b\int_0^b\theta'_{z\alpha\beta}
\le\Bover1{2^lm}\int_0^{2^l(m+1)}\theta'_{z\alpha\beta}
=\Bover{m+1}{m}\gamma$,}

\noindent dan kedua batas mendekati $\gamma$ saat $b\to\infty$.\ \Qed

\medskip

{\bf (c)} Karena $(\alpha,y)\mapsto\theta'_{z\alpha\beta}(y)$ selalu Borel
terukur, setiap fungsi
$\alpha\mapsto\Bover1n\int_0^n\theta'_{z\alpha\beta}$, untuk $n\ge 1$, adalah
Borel terukur (dengan menggabungkan 251M dan 252P), dan $\alpha\mapsto
g(\alpha,y,z):\ooint{0,\infty}\to\Bbb R$ adalah Borel terukur; pada saat yang
sama, karena $0\le\theta'_{z\alpha\beta}(y)\le 1$ untuk semua $\alpha$ dan
$\beta$, $0\le g(\alpha,y,z)\le 1$ untuk setiap $\alpha$, dan
$\tilde\theta_z(y)=\int_1^2\Bover1{\alpha}g(\alpha,y,z)d\alpha$ didefinisikan
dalam $[0,1]$.

\medskip

{\bf (d) } Untuk setiap $y<z$, $\gamma\in\Bbb R$ dan $\alpha>0$,
$g(\alpha,y+\gamma,z+\gamma)=g(\alpha,y,z)$. \Prf\ Cukup untuk
mempertimbangkan kasus $\gamma\ge 0$. Dalam hal ini

$$\eqalign{g(\alpha,y+\gamma,z+\gamma)
&=\lim_{b\to\infty}\Bover1b\int_0^b
  \theta'_{z+\gamma,\alpha,\beta}(y+\gamma)d\beta\cr
&=\lim_{b\to\infty}\Bover1b\int_0^b
  \theta_{\alpha z+\alpha\gamma+\beta}
    (\alpha y+\alpha\gamma+\beta)d\beta\cr
&=\lim_{b\to\infty}\Bover1b\int_{\alpha\gamma}^{b+\alpha\gamma}
  \theta_{\alpha z+\beta}(\alpha y+\beta)d\beta
=\lim_{b\to\infty}\Bover1b\int_{\alpha\gamma}^{b+\alpha\gamma}
  \theta'_{z\alpha\beta}(y)d\beta,\cr}$$

\noindent jadi

$$\eqalign{|g(\alpha,y+\gamma,z+\gamma)-g(\alpha,y,z)|
&=\lim_{b\to\infty}\Bover1b\bigl
  |\int_b^{b+\alpha\gamma}\theta'_{z\alpha\beta}(y)d\beta
  -\int_0^{\alpha\gamma}\theta'_{z\alpha\beta}(y)d\beta\bigr|\cr
&\le\lim_{b\to\infty}\Bover{2\alpha\gamma}{b}=0. \text{\Qed}\cr}$$

Oleh karena itu, setiap kali $y<z$ dan $\gamma\in\Bbb R$,

\Centerline{$\tilde\theta_{z+\gamma}(y+\gamma)
=\int_{0}^{1}\Bover1{\alpha}g(\alpha,y+\gamma,z+\gamma)d\alpha
=\int_{0}^{1}\Bover1{\alpha}g(\alpha,y,z)d\alpha =\tilde\theta_z(y)$.}

\medskip

{\bf (e) } Fakta penting berikutnya yang perlu dicatat adalah bahwa
$\theta_{2z}(2y)$ selalu sama dengan $\theta_z(y)$. \Prf\ Jika $\theta_z(y)\ne
0$, maka (seperti dalam (b) di atas) ada $k$, $m\in\Bbb Z$ sedemikian sehingga

\Centerline{$2^k(m+\Bover12)\le z<2^k(m+1)$, \quad$2^km\le y<2^k(m+\Bover12)$,
\quad$\theta_z(y)=\varhat{\phi}(2^{-k}y-(m+\Bover14))^2$.}

\noindent Dalam hal ini,

\Centerline{$2^{k+1}(m+\Bover12)\le 2z<2^{k+1}(m+1)$, \quad$2^{k+1}m\le
2y<2^{k+1}(m+\Bover12)$,}

\noindent jadi

\Centerline{$\theta_{2z}(2y) =\varhat{\phi}(2^{-k-1}\cdot 2y-(m+\Bover14))^2
=\theta_z(y)$.}

\noindent Sama halnya,

\Centerline{$2^{k-1}(m+\Bover12)\le \Bover12z<2^{k-1}(m+1)$, \quad$2^{k-1}m\le
\Bover12y<2^{k-1}(m+\Bover12)$,}

\noindent jadi

\Centerline{$\theta_{\bover12z}(\Bover12y)
=\varhat{\phi}(2^{-k+1}\cdot\Bover12y-(m+\Bover14))^2 =\theta_z(y)$.}

\noindent Ini menunjukkan bahwa $\theta_{2z}(2y)=\theta_z(y)$ jika salah satu
ruas tidak nol, dan karena itu dalam semua kasus.\ \Qed

Oleh karena itu

\Centerline{$\theta'_{z,2\alpha,2\beta}(y) =\theta_{2\alpha z+2\beta}(2\alpha
y+2\beta) =\theta_{\alpha z+\beta}(\alpha y+\beta)
=\theta'_{z\alpha\beta}(y)$}

\noindent untuk semua $y$, $z$, $\beta\in\Bbb R$ dan semua $\alpha>0$.

\medskip

{\bf (f)} Akibatnya

$$\eqalign{g(2\alpha,y,z)
&=\lim_{b\to\infty}\Bover1b\int_0^b\theta'_{z,2\alpha,\beta}(y)d\beta
=\lim_{b\to\infty}\Bover2b\int_0^{b/2}
  \theta'_{z,2\alpha,2\beta}(y)d\beta\cr
&=\lim_{b\to\infty}\Bover2b\int_0^{b/2}
  \theta'_{z\alpha\beta}(y)d\beta
=\lim_{b\to\infty}\Bover1b\int_0^b\theta'_{z\alpha\beta}(y)d\beta
=g(\alpha,y,z)\cr}$$

\noindent setiap kali $\alpha>0$ dan $y$, $z\in\Bbb R$.

\Centerline{$\int_{\gamma}^{\delta} \Bover1{\alpha}g(\alpha,y,z)d\alpha
=\int_{\gamma}^{\delta}\Bover1{\alpha}g(2\alpha,y,z)d\alpha
=\int_{2\gamma}^{2\delta}\Bover1{\alpha}g(\alpha,y,z)d\alpha$}

\noindent setiap kali $0<\gamma\le\delta$, dan karena itu

\Centerline{$\int_{\gamma}^{2\gamma} \Bover1{\alpha}g(\alpha,y,z)d\alpha
=\int_1^2\Bover1{\alpha}g(\alpha,y,z)d\alpha$}

\noindent untuk setiap $\gamma>0$. \Prf\ Pilih $k\in\Bbb Z$ sehingga
$2^k\le\gamma<2^{k+1}$. Maka

$$\eqalign{\int_{\gamma}^{2\gamma}\Bover1{\alpha}g(\alpha,y,z)d\alpha
&=\int_{2^k}^{2^{k+1}}\Bover1{\alpha}g(\alpha,y,z)d\alpha
  -\int_{2^k}^{\gamma}\Bover1{\alpha}g(\alpha,y,z)d\alpha
  +\int_{2^{k+1}}^{2\gamma}\Bover1{\alpha}g(\alpha,y,z)d\alpha\cr
&=\int_{2^k}^{2^{k+1}}\Bover1{\alpha}g(\alpha,y,z)d\alpha
=\int_1^2\Bover1{\alpha}g(\alpha,y,z)d\alpha.\text{\Qed}\cr}$$

\medskip

{\bf (g)} Sekarang jika $\alpha$, $\gamma>0$ dan $y<z$,

\Centerline{$g(\alpha,\gamma y,\gamma z) =\lim_{b\to\infty}\Bover1b\int_0^b
\theta_{\alpha\gamma z+\beta}(\alpha\gamma y+\beta)d\beta
=g(\alpha\gamma,y,z)$.}

\noindent Jadi jika $\gamma>0$ dan $y<z$,

$$\eqalign{\tilde\theta_{\gamma z}(\gamma y)
&=\int_1^2\Bover1{\alpha}g(\alpha,\gamma y,\gamma z)d\alpha
=\int_1^2\Bover1{\alpha}g(\alpha\gamma,y,z)d\alpha\cr
&=\int_{\gamma}^{2\gamma}\Bover1{\alpha}g(\alpha,y,z)d\alpha
=\int_1^2\Bover1{\alpha}g(\alpha,y,z)d\alpha
=\tilde\theta_z(y).\cr}$$

\noindent Menggabungkan ini dengan (d), kita melihat bahwa jika $y<z$ maka

\Centerline{$\tilde\theta_z(y) =\tilde\theta_{z-y}(0) =\tilde\theta_1(0)$.}

\medskip

{\bf (h) } Masih perlu diperiksa bahwa $\tilde\theta_1(0)$ tidak nol.
Anggap $1\le\alpha<\bover76$ dan terdapat $m\in\Bbb Z$ sedemikian sehingga
$2(m+\bover1{12})\le\beta\le 2(m+\bover5{12})$. Maka
$2(m+\bover12)\le\alpha+\beta<2(m+1)$, sementara
$|\bover12\beta-(m+\bover14)|\le\bover16$, jadi

\Centerline{$\theta_{\alpha+\beta}(\beta)
=\varhat{\phi}(\Bover12\beta-(m+\Bover14))^2 =1$.}

\noindent Ini berarti bahwa, untuk $1\le\alpha<\bover76$,

$$\eqalign{g(\alpha,0,1)
&=\lim_{m\to\infty}\Bover1{2m}\int_0^{2m}
  \theta_{\alpha+\beta}(\beta)d\beta\cr
&\ge\lim_{m\to\infty}\Bover{1}{2m}
  \sum_{j=0}^{m-1}\mu[2(j+\Bover1{12}),2(j+\Bover5{12})]
=\Bover13.\cr}$$

\noindent Jadi

\Centerline{$\tilde\theta_1(0) =\int_1^2\Bover1{\alpha}g(\alpha,0,1)d\alpha
\ge\Bover13\int_1^{7/6}\Bover1{\alpha}d\alpha >0$.}

\noindent Ini menyelesaikan bukti.
}%end of proof of 286R

\leader{286S}{Lema} Misalkan $h$ fungsi uji yang menurun cepat.

(a) Untuk setiap $x\in\Bbb R$,

\Centerline{$(\tilde Ah)(x)
=\liminf_{n\to\infty}\Bover1n\int_1^2\Bover1{\alpha}\int_0^n
(D_{1/\alpha}AM_{\beta}D_{\alpha}h)(x)d\beta d\alpha$}

\noindent didefinisikan dalam $[0,\infty]$, dan $\tilde Ah:\Bbb
R\to[0,\infty]$ adalah Borel terukur.

(b) $\int_F\tilde Ah\le 3C_9\|h\|_2\sqrt{\mu F}$ kapan pun $\mu F<\infty$.

(c) Jika $z\in\Bbb R$, $2\pi|(\varhat
h\times\tilde\theta_z)\varspcheck|\le\tilde Ah$ di setiap titik.

\proof{{\bf (a)} Intinya di sini adalah bahwa fungsi

\Centerline{$(\alpha,\beta,x) \mapsto(D_{1/\alpha}AM_{\beta}D_{\alpha}h)(x):
\ooint{0,\infty}\times\BbbR^2\to[0,\infty]$}

\noindent adalah Borel terukur. \Prf\

$$\eqalign{(D_{1/\alpha}AM_{\beta}D_{\alpha}h)(x)
&=(AM_{\beta}D_{\alpha}h)(\Bover{x}{\alpha})
\cr&=\sup_{z\in\Bbb R}
  |2\pi((M_{\beta}D_{\alpha}h)\varsphat\times\theta_z)\varspcheck
     (\Bover{x}{\alpha})|
\cr&=\Bover{2\pi}{\alpha}\sup_{z\in\Bbb R}
  |(S_{-\beta}D_{1/\alpha}\varhat h\times\theta_z)\varspcheck
     (\Bover{x}{\alpha})|.
\cr}$$
  
\noindent Sekarang, untuk setiap $z\in\Bbb R$,

\Centerline{$(S_{-\beta}D_{1/\alpha}\varhat h\times\theta_z)\varspcheck
(\bover{x}{\alpha}) =\Bover1{\sqrt{2\pi}}\int_{-\infty}^{\infty}e^{ixy/\alpha}
\varhat h(\bover{y-\beta}{\alpha})\theta_z(y)dy$.}

\noindent Kita tahu bahwa $\varhat h$ merupakan fungsi uji yang menurun cepat,
sehingga ada $\gamma\ge 0$ sedemikian sehingga $|\varhat
h(t)|\le\Bover{\gamma}{1+t^2}$ untuk setiap $t\in\Bbb R$. Ini berarti bahwa jika
$\alpha>0$ dan $\beta\in\Bbb R$ dan $\sequencen{\alpha_n}$,
$\sequencen{\beta_n}$ adalah barisan dalam $\ocint{0,2\alpha}$ dan
$[\beta-1,\beta+1]$, bertepatan dengan $\alpha$, $\beta$ masing-masing, dan
kita menetapkan $g(t) =\sup_{n\in\Bbb N}|\varhat
h(\Bover{t-\beta_n}{\alpha_n})\theta_z(t)|$, maka

$$\eqalign{g(t)
&\le\Bover{4\gamma\alpha^2}{(|t|-|\beta|-1)^2}\text{jika}|t|\ge|\beta|+2,
  \cr
&\le\gamma\text{lain},\cr}$$

\noindent dan $g$ terintegralkan. (Ingat bahwa $0\le\theta_z(y)\le 1$
untuk setiap $y$, seperti dicatat dalam 286Oa.) Jadi Teorema Konvergensi
Terdominasi Lebesgue menyatakan bahwa jika
$\sequencen{\alpha_n}\to\alpha$ dan $\sequencen{\beta_n}\to\beta$ dan
$\sequencen{x_n}\to x$, maka

\Centerline{$\int_{-\infty}^{\infty}e^{ix_ny/\alpha_n} \varhat
h(\bover{y-\beta_n}{\alpha_n})\theta_z(y)dy
\to\int_{-\infty}^{\infty}e^{ixy/\alpha} \varhat
h(\bover{y-\beta}{\alpha})\theta_z(y)dy$.}
   
\noindent Jadi $(\alpha,\beta,x) \mapsto(S_{-\beta}D_{1/\alpha}\varhat
h\times\theta_z)\varspcheck (\bover{x}{\alpha})$ kontinu. Ini berlaku
untuk setiap $z\in\Bbb R$. Oleh karena itu

\Centerline{$(\alpha,\beta,x)\mapsto\sup_{z\in\Bbb R}
|(S_{-\beta}D_{1/\alpha}\varhat h\times\theta_z)\varspcheck
(\Bover{x}{\alpha})|$}

\noindent dan $(\alpha,\beta,x) \mapsto(D_{1/\alpha}AM_{\beta}D_{\alpha}h)(x)$
semikontinu bawah, sehingga terukur Borel, sekali lagi menurut 256Ma
lagi.\ \Qed

Oleh karena itu, integral

\Centerline{$\int_1^2\Bover1{\alpha}\int_0^n
(D_{1/\alpha}AM_{\beta}D_{\alpha}h)(x)d\beta d\alpha$}

\noindent terdefinisi dalam $[0,\infty]$ dan merupakan fungsi terukur Borel
dari $x$ (sekali lagi 252P), sehingga $\tilde Af$ terukur Borel.

\medskip

{\bf (b)} Untuk setiap $n\in\Bbb N$,

$$\eqalignno{\int_F\Bover1n\int_1^2\Bover1\alpha\int_0^n
&(D_{1/\alpha}AM_{\beta}D_{\alpha}h)(x)d\beta d\alpha dx\cr
&=\Bover1n\int_1^2\Bover1\alpha\int_0^n\!\int_F
  (D_{1/\alpha}AM_{\beta}D_{\alpha}h)(x)dx d\beta d\alpha\cr
\displaycause{menurut teorema Fubini, 252H}
&=\Bover1n\int_1^2\!\int_0^n\!\int_F\Bover1\alpha
  (AM_{\beta}D_{\alpha}h)(\Bover{x}{\alpha})dx d\beta d\alpha\cr
&=\Bover1n\int_1^2\!\int_0^n\int_{\alpha^{-1}F}
  (AM_{\beta}D_{\alpha}h)(x)dx d\beta d\alpha\cr
&\le\Bover1n\int_1^2\!\int_0^n 4C_9
  \|M_{\beta}D_{\alpha}h\|_2\sqrt{\mu(\alpha^{-1}F)} d\beta d\alpha\cr
\displaycause{286P}
&=4C_9\cdot\Bover1n\int_1^2\!\int_0^n
  \Bover1{\sqrt{\alpha}}\|h\|_2\cdot\Bover1{\sqrt{\alpha}}\sqrt{\mu F}
    d\beta d\alpha\cr
&=4C_9\|h\|_2\sqrt{\mu F}\cdot\Bover1n\int_1^2\Bover1{\alpha}\int_0^n
 d\beta d\alpha\cr
&=4C_9\|h\|_2\sqrt{\mu F}\ln 2
\le 3C_9\|h\|_2\sqrt{\mu F}.\cr}$$

\noindent Jadi

$$\eqalignno{\int_F\tilde Ah
&=\int_F\liminf_{n\to\infty}\Bover1n\int_1^2\Bover1\alpha\int_0^n
  (D_{1/\alpha}AM_{\beta}D_{\alpha}h)(x)d\beta d\alpha dx\cr
&\le\liminf_{n\to\infty}\int_F\Bover1n\int_1^2\Bover1\alpha\int_0^n
  (D_{1/\alpha}AM_{\beta}D_{\alpha}h)(x)d\beta d\alpha dx\cr
\displaycause{menurut lema Fatou}
&\le 3C_9\|h\|_2\sqrt{\mu F}.\cr}$$

\medskip

{\bf (c)} Untuk setiap $x\in\Bbb R$,

\Centerline{$\int_{-\infty}^{\infty} |\varhat h(y)|\int_1^2\Bover1{\alpha}
\bigl(\sup_{n\in\Bbb N}\Bover1n\int_0^n
\theta'_{z\alpha\beta}(y)d\beta\bigr)d\alpha dy \le\ln
2\cdot\int_{-\infty}^{\infty}|\varhat h|$}

\noindent adalah hingga. Jadi

$$\eqalignno{(\varhat h\times\tilde\theta_z)\varspcheck(x)
&=\Bover1{\sqrt{2\pi}}\int_{-\infty}^{\infty}
  e^{ixy}\varhat h(y)\tilde\theta_z(y)dy\cr
&=\Bover1{\sqrt{2\pi}}\int_{-\infty}^{\infty}
  e^{ixy}\varhat h(y)\int_1^2\Bover1{\alpha}\lim_{n\to\infty}
  \Bover1n\int_0^n\theta'_{z\alpha\beta}(y)d\beta d\alpha dy\cr
&=\Bover1{\sqrt{2\pi}}\lim_{n\to\infty}\int_{-\infty}^{\infty}
  e^{ixy}\varhat h(y)\int_1^2\Bover1{\alpha n}
  \int_0^n\theta'_{z\alpha\beta}(y)d\beta d\alpha dy\cr
\displaycause{menurut Teorema Konvergensi Terdominasi Lebesgue}
&=\Bover1{\sqrt{2\pi}}\lim_{n\to\infty}\int_1^2\Bover1{\alpha n}
  \int_0^n\int_{-\infty}^{\infty}
  e^{ixy}\varhat h(y)\theta'_{z\alpha\beta}(y)dy d\beta d\alpha\cr
\displaycause{menurut teorema Fubini}
&=\lim_{n\to\infty}\int_1^2\Bover1{\alpha n}
  \int_0^n(\varhat h\times\theta'_{z\alpha\beta})\varspcheck(x)
  d\beta d\alpha,\cr}$$

\noindent dan

$$\eqalignno{2\pi|(\varhat h\times\tilde\theta_z)\varspcheck(x)|
&=2\pi\bigl|\lim_{n\to\infty}\int_1^2\Bover1{\alpha n}
  \int_0^n(\varhat h\times\theta'_{z\alpha\beta})\varspcheck(x)
  d\beta d\alpha\bigr|\cr
&\le 2\pi\liminf_{n\to\infty}\int_1^2\Bover1{\alpha n}
  \int_0^n|(\varhat h\times\theta'_{z\alpha\beta})\varspcheck(x)|
  d\beta d\alpha\cr
&\le\liminf_{n\to\infty}\int_1^2\Bover1{\alpha n}
  \int_0^n(D_{1/\alpha}AM_{\beta}D_{\alpha}h)(x)
  d\beta d\alpha\cr
\displaycause{286Qb}
&=(\tilde Ah)(x).\cr}$$
}%end of proof of 286S

\leader{286T}{Lema} Tetapkan $C_{10}=3C_9/\pi\tilde\theta_1(0)$. Untuk $f\in\eusm
L^2_{\Bbb C}$, definisikan $\hat Af:\Bbb R\to[0,\infty]$ dengan menetapkan

\Centerline{$(\hat Af)(y) =\sup_{a\le
b}\Bover1{\sqrt{2\pi}}|\int_a^be^{-ixy}f(x)dx|$}

\noindent untuk setiap $y\in\Bbb R$. Maka $\int_F\hat Af\le
C_{10}\|f\|_2\sqrt{\mu F}$ setiap kali $\mu F<\infty$.

\proof{{\bf (a)} Seperti biasa, langkah pertama adalah memastikan
bahwa $\hat Af$ adalah terukur. \Prf\ Untuk $a\le b$,
$y\mapsto|\Bover1{\sqrt{2\pi}}\int_a^be^{-ixy}f(x)dx|$ adalah kontinu (dengan
283Cf, karena $f\times\chi[a,b]$ adalah terintegralkan), sehingga $\hat Af$
semikontinu bawah, dan karena itu terukur Borel (sekali lagi 256Ma
sekali lagi). \ \Qed

\medskip

{\bf (b)} Misalkan $h$ fungsi uji yang menurun cepat. Maka

\Centerline{$(\hat Ah)(y) \le\Bover1{\pi\tilde\theta_1(0)}(\tilde A\varcheck
h)(-y)$}

\noindent untuk setiap $y\in\Bbb R$. \Prf\ Jika $a\in\Bbb R$ maka

$$\eqalignno{\Bover1{\sqrt{2\pi}}|\int_{-\infty}^ae^{-ixy}h(x)dx|
&=\Bover1{\tilde\theta_1(0)\sqrt{2\pi}}|\int_{-\infty}^{\infty}
  e^{-ixy}\tilde\theta_a(x)h(x)dx|\cr
\displaycause{286R}
&=\Bover1{\tilde\theta_1(0)}|(h\times\tilde\theta_a)\varspcheck(-y)|
=\Bover1{\tilde\theta_1(0)}
  |(\varcheck h\varhat{\phantom{h}}
    \times\tilde\theta_a)\varspcheck(-y)|\cr
\displaycause{sekali lagi menurut 284C}
&\le\Bover1{2\pi\tilde\theta_1(0)}(\tilde A\varcheck{h})(-y)\cr}$$

\noindent (286Sc). Jadi jika $a\le b$ di dalam $\Bbb R$,

\Centerline{$\Bover1{\sqrt{2\pi}}|\int_{a}^{b}e^{-ixy}h(x)dx|
\le\Bover1{\pi\tilde\theta_1(0)}(\tilde A\varcheck{h})(-y)$;}

\noindent dengan mengambil supremum atas $a$ dan $b$, diperoleh hasilnya.\ \Qed

Ini berarti bahwa

$$\eqalignno{\int_F\hat Ah
&\le\Bover1{\pi\tilde\theta_1(0)}\int_{-F}\tilde A\varcheck h
\le\Bover3{\pi\tilde\theta_1(0)}C_9\|h\|_2\sqrt{\mu(-F)}\cr
\displaycause{286Sb, 284Oa}
&=C_{10}\|h\|_2\sqrt{\mu F}.\cr}$$

\medskip

{\bf (c)} Untuk $f$ terintegralkan-kuadrat sembarang, ambil $\epsilon>0$
dan $n\in\Bbb N$ sembarang. Tetapkan

\Centerline{$(\hat A_nf)(y) =\sup_{-n\le a\le b\le
n}\Bover1{\sqrt{2\pi}}|\int_a^be^{-ixy}f(x)dx|$}

\noindent untuk setiap $y\in\Bbb R$. Misalkan $h$ fungsi uji yang menurun
cepat sedemikian sehingga $\|f-h\|_2\le\epsilon$ (284N). Maka

\Centerline{$\hat Ah \ge\hat A_nh \ge\hat
A_nf-\Bover{\sqrt{2n}}{\sqrt{2\pi}}\epsilon$}

\noindent (menggunakan ketidaksetaraan Cauchy), jadi

\Centerline{$\int_F\hat A_nf \le\int_F\hat Ah+\sqrt{\bover{n}{\pi}}\epsilon\mu
F \le C_{10}(\|f\|_2+\epsilon)\sqrt{\mu F} +\sqrt{\bover{n}{\pi}}\epsilon\mu
F$.}

\noindent Karena $\epsilon$ sembarang, $\int_F\hat A_nf\le C_{10}\|f\|_2\sqrt{\mu F}$; dengan membiarkan $n\to\infty$, kita memperoleh $\int_F\hat Af\le C_{10}\|f\|_2\sqrt{\mu F}$. }%end of proof of 286T

\vleader{72pt}{286U}{Teorema} Jika $f\in\eusm L^2_{\Bbb C}$ maka

\Centerline{$g(y) =\lim_{a\to-\infty,b\to\infty}
\Bover1{\sqrt{2\pi}}\int_a^be^{-ixy}f(x)dx$}

\noindent terdefinisi dalam $\Bbb C$ untuk hampir setiap $y\in\Bbb R$, dan
$g$ mewakili transformasi Fourier dari $f$.

\proof{{\bf (a)} Untuk $n\in\Bbb N$, $y\in\Bbb R$, tetapkan

\Centerline{$\gamma_n(y)=\sup_{a\le -n,b\ge n}\Bover1{\sqrt{2\pi}}
\bigl|\int_a^be^{-ixy}f(x)dx-\int_{-n}^ne^{-ixy}f(x)dx\bigr|$.}

\noindent Maka $g(y)$ terdefinisi setiap kali $\inf_{n\in\Bbb
N}\gamma_n(y)=0$. \Prf\ Jika $\inf_{n\in\Bbb N}\gamma_n(y)=0$ dan
$\epsilon>0$, pilih $m\in\Bbb N$ sehingga $\gamma_m(y)\le\bover12\epsilon$;
maka $\Bover1{\sqrt{2\pi}}
|\int_a^be^{-ixy}f(x)dx-\int_{-n}^ne^{-ixy}f(x)dx|\le\epsilon$ setiap kali
$n\ge m$, $a\le-n$, dan $b\ge n$. Ini berarti, pertama, bahwa
$\sequencen{\int_{-n}^ne^{-ixy}f(x)dx}$ merupakan barisan Cauchy, sehingga
memiliki limit, katakanlah $\zeta$, dan, kedua, bahwa
$\zeta=\lim_{a\to-\infty,b\to\infty}\int_a^be^{-ixy}f(x)dx$, sehingga
$g(y)=\Bover{\zeta}{\sqrt{2\pi}}$ didefinisikan.\ \Qed

Setiap $\gamma_n$ juga semikontinu bawah (bdk.\ bagian (a) dari
bukti dari 286T).

\medskip

{\bf (b)} \Quer\ Misalkan, jika mungkin, bahwa $\{y:\inf_{n\in\Bbb
N}\gamma_n(y)>0\}$ tidak dapat diabaikan.

\Centerline{$\lim_{m\to\infty}\mu\{y:|y|\le m,\, \inf_{n\in\Bbb
N}\gamma_n(y)\ge\Bover1m\} >0$,}

\noindent jadi ada $\epsilon>0$ sedemikian sehingga

\Centerline{$F=\{y:|y|\le\Bover1{\epsilon},\, \inf_{n\in\Bbb
N}\gamma_n(y)\ge\epsilon\}$}

\noindent berukuran lebih besar dari $\epsilon$. Pilih $n\in\Bbb N$
sedemikian sehingga

\Centerline{$4C^2_{10}
(\int_{-\infty}^{\infty}|f(x)|^2dx-\int_{-n}^n|f(x)|^2dx) \le\epsilon^3$,}

\noindent dan tetapkan $f_1=f-f\times\chi[-n,n]$; maka
$2C_{10}\|f_1\|_2\le\epsilon^{3/2}$.

Kita mempunyai

$$\eqalign{\gamma_n(y)
&=\sup_{a\le -n,b\ge n}\Bover1{\sqrt{2\pi}}
  \bigl|\int_a^be^{-ixy}f_1(x)dx-\int_{-n}^ne^{-ixy}f_1(x)dx\bigr|\cr
&\le 2\sup_{a\le b}\Bover1{\sqrt{2\pi}}
  |\int_a^be^{-ixy}f_1(x)dx|
\le 2(\hat Af_1)(y),\cr}$$

\noindent sehingga

$$\eqalignno{\epsilon\mu F
&\le\int_F\gamma_n
\le 2\int_F\hat Af_1
\le 2C_{10}\|f_1\|_2\sqrt{\mu F}\cr
\displaycause{286T}
&\le\epsilon^{3/2}\sqrt{\mu F}\cr}$$

\noindent dan $\mu F\le\epsilon$; tetapi kita memilih $\epsilon$ sehingga $\mu
F$ lebih besar dari $\epsilon$.\ \Bang

\medskip

{\bf (c)} Jadi $g(y)$ terdefinisi untuk hampir setiap $y\in\Bbb R$.
Sekarang $g$ mewakili transformasi Fourier dari $f$. \Prf\ Misalkan $h$
fungsi uji yang menurun cepat. Pembatasan $\hat Af$ pada himpunan tempat
nilainya berhingga merupakan fungsi bertemper, menurut 286D. Jadi
$\int_{-\infty}^{\infty}(\hat Af)\times|h|$ adalah hingga, oleh 284F.
Sekarang

$$\eqalignno{\int_{-\infty}^{\infty}g\times h
&=\Bover1{\sqrt{2\pi}}\int_{-\infty}^{\infty}
  \bigl(\lim_{n\to\infty}\int_{-n}^ne^{-ixy}f(x)dx\bigr)h(y)dy\cr
&=\Bover1{\sqrt{2\pi}}\lim_{n\to\infty}\int_{-\infty}^{\infty}
  \int_{-n}^ne^{-ixy}f(x)h(y)dxdy\cr
\displaycause{karena
$\Bover1{\sqrt{2\pi}}|\int_{-n}^ne^{-ixy}f(x)dx|\le\hat Af(y)$ untuk setiap
$n$ dan $y$, sehingga Teorema Konvergensi Terdominasi Lebesgue dapat digunakan}
&=\Bover1{\sqrt{2\pi}}\lim_{n\to\infty}
  \int_{-n}^n\int_{-\infty}^{\infty}e^{-ixy}f(x)h(y)dydx\cr
\displaycause{karena $\int_{-\infty}^{\infty}\int_{-n}^n|f(x)h(y)|dxdy$
berhingga untuk setiap $n$}
&=\lim_{n\to\infty}
  \int_{-n}^nf\times\varhat{h}
=\int_{-\infty}^{\infty}f\times\varhat{h}\cr}$$

\noindent karena $f\times\varhat{h}$ terintegralkan. Karena $h$ sembarang, $g$ mewakili transformasi Fourier dari $f$.\ \Qed }%end of proof of 286U

\leader{286V}{Teorema} Untuk setiap fungsi bernilai kompleks yang
terintegralkan-kuadrat pada $\ocint{-\pi,\pi}$, barisan jumlah parsial
Fouriernya konvergen kepadanya hampir di mana-mana.

\proof{ Misalkan $f\in\eusm L^2_{\Bbb C}(\mu_{\ocint{-\pi,\pi}})$. Tetapkan
$f_1(x)=f(x)$ untuk $x\in\dom f$, $0$ untuk $x\in\Bbb
R\setminus\ocint{-\pi,\pi}$; maka $f_1\in\eusm L^2_{\Bbb C}(\mu)$. Misalkan
$g\in\eusm L^2_{\Bbb C}(\mu)$ mewakili transformasi Fourier invers dari $f_1$
(284O). Menurut 286U,
$f_2(x)=\lim_{a\to\infty}\Bover1{\sqrt{2\pi}}\int_{-a}^ae^{-ixy}g(y)dy$
terdefinisi untuk hampir setiap $x$, dan $f_2$ mewakili transformasi Fourier
dari $g$, sehingga sama dengan $f_1$ hampir di mana-mana (284Ib).

Sekarang, untuk setiap $a\ge 0$, $x\in\Bbb R$,

$$\eqalignno{\int_{-a}^ae^{-ixy}g(y)dy
&=\innerprod{g}{h_{ax}}\cr
\displaycause{dengan $h_{ax}(y)=e^{ixy}$ jika $|y|\le a$, dan $0$ selainnya}
&=\innerprod{f_2}{\varhat{h}_{ax}}\cr
\displaycause{284Ob}
&=\Bover1{\sqrt{2\pi}}\int_{-\infty}^{\infty}f_2(t)
  \overline{\int_{-\infty}^{\infty}e^{-ity}h_{ax}(y)dy}\,dt\cr
&=\Bover2{\sqrt{2\pi}}\int_{-\infty}^{\infty}
  \Bover{\sin(x-t)a}{x-t}f_2(t)dt
=\Bover2{\sqrt{2\pi}}\int_{-\pi}^{\pi}
  \Bover{\sin(x-t)a}{x-t}f(t)dt.\cr}$$

\noindent Jadi

\Centerline{$f(x)=f_2(x) =\lim_{a\to\infty}\Bover1{\pi}\int_{-\pi}^{\pi}
\Bover{\sin(x-t)a}{x-t}f(t)dt$}

\noindent untuk hampir setiap $x\in\ocint{-\pi,\pi}$.

Di sisi lain, dengan menuliskan $\sequencen{s_n}$ untuk barisan jumlah parsial
Fourier dari $f$, untuk setiap $x\in\ooint{-\pi,\pi}$ kita mempunyai

\Centerline{$s_n(x) =\Bover1{2\pi}\int_{-\pi}^{\pi}f(t) \Bover{\sin(n+{1\over
2}\pushbottom{3.5pt})(x-t)} {\sin{1\over 2}(x-t)}dt$}

\noindent untuk setiap $n$, dengan 282Da. Sekarang

$$\eqalign{\Bover1{2\pi}\int_{-\pi}^{\pi}f(t)
  &\Bover{\sin(n+{1\over 2}\pushbottom{3.5pt})(x-t)}
  {\sin{1\over 2}(x-t)}dt
-\Bover1{\pi}\int_{-\pi}^{\pi}f(t)
  \Bover{\sin(n+{1\over 2}\pushbottom{3.5pt})(x-t)}
  {x-t}dt\cr
&=\Bover1{\pi}\int_{-\pi}^{\pi}f(t)
  \bigl(\Bover{\sin(n+{1\over 2}\pushbottom{3.5pt})(x-t)}
    {2\sin{1\over 2}(x-t)}
  -\Bover{\sin(n+\bover12\pushbottom{3.5pt})(x-t)}{x-t}\bigr)dt\cr
&=\Bover1{\pi}\int_{x-\pi}^{x+\pi}\bigl
  (\Bover1{2\sin{1\over 2}t}-\Bover1{t}\bigr)
  f(x-t)\sin(n+\hbox{$\bover12$})t\,dt.\cr}$$

\noindent Tetapi jika kita meninjau fungsi

$$\eqalign{p_x(t)
&=\bigl(\Bover1{2\sin\bover12t}-\Bover1t\bigr)f(x-t)
  \text{ jika }x-\pi<t<x+\pi\text{ dan }t\ne 0,\cr
&=0\text{ selainnya},\cr}$$

\noindent $p_x$ terintegralkan, karena $f$ terintegralkan pada
$\ocint{-\pi,\pi}$ dan $\lim_{t\to 0}\Bover1{2\sin\bover12t}-\Bover1t=0$, jadi
$\sup_{t\ne 0,x-\pi\le t\le x+\pi}|\Bover1{2\sin\bover12t}-\Bover1t|$ adalah
berhingga. (Di sini kita menggunakan $|x|<\pi$.) Jadi

$$\eqalign{\lim_{n\to\infty}s_n(x)-\Bover1{\pi}\int_{-\pi}^{\pi}f(t)
  \Bover{\sin(n+{1\over 2}\pushbottom{3.5pt})(x-t)}
  {x-t}dt
=\lim_{n\to\infty}\int_{-\infty}^{\infty}
       p_x(t)\sin(n+\hbox{$\bover12$})t\,dt
=0\cr}$$

\noindent menurut lema Riemann--Lebesgue (282Fb). Ini berarti bahwa $\lim_{n\to\infty}s_n(x)=f(x)$ untuk setiap $x\in\ooint{-\pi,\pi}$ sedemikian sehingga $f(x)=\lim_{a\to\infty}\Bover1{\pi}\int_{-\pi}^{\pi}
  \Bover{\sin(x-t)a}{x-t}f(t)dt$, yang hampir setiap $x\in\ocint{-\pi,\pi}$.
}%end of proof of 286V

\cmmnt{ \vleader{108pt}{286W}{Glosarium} Notasi khusus berikut digunakan dalam
lebih dari satu paragraf dari bagian ini:

\def\vta{$\mu$ untuk ukuran Lebesgue pada $\Bbb R$.} \def\vtb{286A: $f^*$.}
\def\vtc{286C: $S_{\alpha}f$, $M_{\alpha}f$, $D_{\alpha}f$.} \def\vtd{286Ea:
$\Cal I$, $Q$, $I_{\sigma}$, $J_{\sigma}$, $k_{\sigma}$, $x_{\sigma}$,
$y_{\sigma}$, $J^l_{\sigma}$, $J^r_{\sigma}$, $y^l_{\sigma}$.} \def\vte{286Eb:
$\phi$, $\phi_{\sigma}$, $\innerprod{f}{g}$.} \def\vtf{286Ec: $w$,
$w_{\sigma}$.} \def\vtg{286F: $\le$, $R^+$, $T_{\tau}$.} \def\vth{286G: $C_1$,
$C_2$, $C_3$, $C_4$.} \def\vti{286H: $\mass$, $\Delta_f$, $\energy$.}
\def\vtj{286J: $C_5$.} \def\vtk{286K: $C_6$.} \def\vtl{286L: $C_7$.}
\def\vtm{286M: $C_8$.} \def\vtn{286N: $C_9$.} \def\vto{286O: $\theta_z$, $\Cal
F$.} \def\vtp{286P: $Ah$.} \def\vtq{286Q: $\theta'_{z\alpha\beta}$.}
\def\vtr{286R: $\tilde\theta_z$.} \def\vts{286S: $\tilde Ah$.} \def\vtt{286T:
$C_{10}$, $\hat Af$.} \sparewidth=\pagewidth \advance\sparewidth by -370pt
\divide\sparewidth by 4

\medskip

\halign{\hskip\sparewidth#\hfil&\hskip\sparewidth#\hfil
&\hskip\sparewidth#\hfil\cr
\vta&\vth&\vto\cr
\vtb&\vti&\vtp\cr
\vtc&\vtj&\vtq\cr
\vtd&\vtk&\vtr\cr
\vte&\vtl&\vts\cr
\vtf&\vtm&\vtt\cr
\vtg&\vtn\cr}
}%end of comment

\exercises{\leader{286X}{Latihan dasar (a)}
%\spheader 286Xa
Gunakan 284Oa dan 284Xg untuk mempersingkat bagian (c) dari bukti 286U.
%286U

\spheader 286Xb Tunjukkan bahwa jika $\sequence{k}{c_k}$ adalah barisan
bilangan kompleks sedemikian sehingga $\sum_{k=0}^{\infty}|c_k|^2$ berhingga, maka
$\sum_{k=0}^{\infty}c_ke^{ikx}$ terdefinisi dalam $\Bbb C$ untuk hampir setiap
$x\in\Bbb R$.
%286V

\leader{286Y}{Latihan lanjutan (a)}
%\spheader 286Ya
Tunjukkan bahwa jika $f$ adalah fungsi terintegralkan-kuadrat pada
$\BbbR^r$, dengan $r\ge 2$, maka

\Centerline{$g(y)
=\Bover1{(\sqrt{2\pi})^r}
 \lim_{\alpha_1,\ldots,\alpha_r\to-\infty,
  \beta_1,\ldots,\beta_r\to\infty}
 \int_{a}^{b}e^{-iy\dotproduct x}f(x)dx$}

\noindent terdefinisi dalam $\Bbb C$ untuk hampir setiap $y\in\BbbR^r$, dan
bahwa $g$ merepresentasikan transformasi Fourier dari $f$.
}%end of exercises

\endnotes{
\Notesheader{286} Ini bukanlah bagian terpanjang dalam keseluruhan
risalah ini, tetapi bagian ini jauh lebih panjang daripada bagian lain mana pun
dalam jilid ini, dan tiga puluh halaman yang dipenuhi subskrip dan superskrip tentu menguji daya tahan
pembaca yang paling bersemangat sekalipun. Anda akan mudah memahami mengapa
Teorema Carleson biasanya tidak disajikan pada tingkat ini. Namun dalam
buku ini saya berusaha menyajikan bukti-bukti lengkap untuk teorema-teorema utama;
dalam jilid-jilid selanjutnya sebagaimana direncanakan saat ini, tidak ada
tempat alami bagi Teorema Carleson; dan teorema itu masih (meski nyaris) dapat
dijangkau pada tahap ini.
Karena itu, saya menyediakan ruang untuk membuktikannya di sini.

Bukti di sini terbagi secara alami menjadi dua bagian: bagian `kombinatorik'
dalam 286E-286M, sampai dengan Lema Lacey--Thiele, disusul bagian
`analitik' dalam 286N-286V, yang menggunakan proses perataan

\Centerline{$\int_1^2\Bover1{\alpha}\lim_{b\to\infty}
\Bover1b\int_0^b\ldots d\beta d\alpha$}

\noindent untuk mengubah fungsi-fungsi $\theta_z$ yang koheren secara geometris,
tetapi tidak reguler secara analitis, menjadi fungsi-fungsi karakteristik
$\Bover1{\tilde\theta_1(0)}\tilde\theta_z$. Dari sudut pandang
analisis Fourier biasa, bagian kedua ini pada dasarnya rutin; ada banyak
jalur yang dapat kita tempuh, dan kita hanya perlu mengambil langkah pencegahan
biasa terhadap operasi yang tidak sah.
\footnote{Pada titik ini saya patut mengakui bahwa saya melakukan
kekeliruan besar dalam edisi 2001 jilid ini dan tidak menyadari kesalahan
tersebut sampai A.Derighetti menunjukkannya kepada saya pada akhir 2013.
Saya berharap versi yang disajikan di sini pada pokoknya benar.}

Carleson ({\smc Carleson 66}) menyatakan teoremanya dalam bentuk deret Fourier
seperti pada 286V; tetapi telah lama dipahami bahwa bentuk ini
ekuivalen dari segi nilai kebenaran
dengan versi transformasi Fourier dalam 286U. Tentu ada banyak cara
untuk memperluas teorema tersebut. Secara khusus, terdapat hasil-hasil
yang bersesuaian untuk fungsi dalam $\eusm L^p$ bagi setiap $p>1$, bahkan untuk
fungsi $f$ sedemikian sehingga $f\times\ln(1+|f|)\times\ln\ln\ln(16+|f|)$
terintegralkan ({\smc Antonov 96}). Metode-metode di sini tampaknya tidak
menjangkau sejauh itu. Perlu saya tambahkan bahwa jika kita mendefinisikan
$\hat Af$ seperti dalam 286T, maka untuk setiap $p>1$ terdapat konstanta
$C$ sedemikian sehingga $\|\hat Af\|_p\le C\|f\|_p$ untuk setiap
$f\in\eusm L^p_{\Bbb C}$
({\smc Hunt 67}, {\smc Mozzochi 71}, {\smc J{\o}rsboe \& Mejlbro 82},
{\smc Arias de Reyna 02}, {\smc Lacey 05}).

Perhatikan bahwa pokok Teorema Carleson, dalam kedua bentuknya, adalah bahwa
kita mengambil limit-limit khusus. Dalam kedua rumus berikut,

\Centerline{$\varhatf(y)
=\Bover1{\sqrt{2\pi}}\lim_{a\to-\infty,b\to\infty}
\int_a^be^{-ixy}f(x)dx$,}

\Centerline{$f(x)=\lim_{n\to\infty}\sum_{-n}^nc_ke^{ikx}$,}

\ifdim\pagewidth>467pt\tenrmstretch{3pt}\fi
\noindent yang berlaku hampir di mana-mana untuk fungsi terintegralkan-kuadrat $f$,
kita tidak mengambil integral biasa
$\int_{-\infty}^{\infty}e^{-ixy}f(x)dx$ ataupun jumlah tak bersyarat
$\sum_{k\in\Bbb Z}c_ke^{ikx}$. Jika $f$ tidak terintegralkan, atau jika
$\sum_{k=-\infty}^{\infty}|c_k|$ tak hingga, bentuk-bentuk ini tidak akan terdefinisi
bahkan pada satu titik pun.
% Weyl's equidistribution theorem
Teorema Carleson bermakna hanya karena kita mempunyai preferensi alami
terhadap jenis-jenis tertentu dari integral tak wajar dan jumlah bersyarat.
Jadi, ketika pada Bab 44 Jilid 4 kita kembali ke analisis Fourier pada grup
topologis umum, sama sekali tidak akan ada bahasa untuk menyatakan teorema
tersebut. Walaupun versi-versinya telah dibuktikan untuk grup lain (misalnya,
{\smc Schipp 78}),
versi-versi itu pasti bergantung pada struktur di luar gagasan sederhana
`grup topologis abelian Hausdorff kompak lokal'. Bahkan dalam
$\BbbR^2$, sejauh yang saya pahami, masih belum diketahui apakah
\tenrmstretch{1.67pt}

\Centerline{$\lim_{a\to\infty}\Bover1{2\pi}\int_{B(\tbf{0},a)}
e^{-iy\dotproduct x}f(x)dx$}

\noindent akan terdefinisi h.d.m.\ untuk setiap fungsi terintegralkan-kuadrat $f$,
jika kita menggunakan bola Euklides biasa $B(\tbf{0},a)$ sebagai pengganti
persegi panjang dalam 286Ya.
}%end of notes

\frnewpage
